
\documentclass{ieeeaccess}

\usepackage[utf8]{inputenc}
\usepackage{cite}
\usepackage{amsmath,amssymb,amsfonts}
\usepackage{amsthm}
\usepackage{float}
\usepackage{graphicx}

\floatstyle{ruled}
\newfloat{algorithm}{!t}{loa}
\floatname{algorithm}{Algorithm}
\usepackage{textcomp}
\usepackage{xcolor}
\usepackage{booktabs}
\usepackage{array}
\usepackage{url}

\newtheorem{theorem}{Theorem}
\newtheorem{proposition}{Proposition}
\newtheorem{definition}{Definition}

\def\BibTeX{{\rm B\kern-.05em{\sc i\kern-.025em b}\kern-.08em
    T\kern-.1667em\lower.7ex\hbox{E}\kern-.125emX}}

\begin{document}

\history{Date of publication xxxx 00, 0000, date of current version xxxx 00, 0000.}
\doi{10.1109/ACCESS.2026.DOI}

\title{Congestion-Tolerant and Auditable Trust Management for the
Internet of Digital Twins}

\author{\uppercase{Abdulelah Alwabel}\authorrefmark{1}, \IEEEmembership{Member, IEEE}
}

\address[1]{Department of Computer Sciences,
    College of Computer Engineering and Sciences,
    Prince Sattam bin Abdulaziz University,
    Al-Kharj, Saudi Arabia}

\tfootnote{This work was supported by Prince Sattam bin Abdulaziz
University under Project PSAU/2026/R/1448.}

\markboth
{A. Alwabel: Congestion-Tolerant and Auditable Trust Management for IoDT}
{A. Alwabel: Congestion-Tolerant and Auditable Trust Management for IoDT}

\corresp{Corresponding author: Abdulelah Alwabel (e-mail: a.alwabel@psau.edu.sa)}

\begin{abstract}
The Internet of Digital Twins (IoDT) trust model mirrors every IoT node
with a cloud-resident twin and isolates a node when any of four binary
factors (energy, packet delivery ratio, traceability and twin
notification) fails.  This rule is fragile under congestion, trusts a
single twin, and keeps verdicts in a mutable database.  We propose
three extensions that keep the factor set: a weighted score smoothed by
an exponential moving average with a per-factor persistence check;
quorum voting over signed votes of replicated twins; and a signed
Merkle audit log kept off the decision path.  In simulation, any rule
requiring two failures removes most false isolations at low
congestion, but only the smoothed score keeps them low as congestion
rises: at $4\%$ per-factor congestion it falsely isolated $1.1\%$ of
honest nodes within 16 rounds, against $93\%$ for the AND and
$14$--$26\%$ for two-failure rules, at the cost of one extra round.
On two public packet-delivery traces it still had the lowest false
isolation in every setting, but real losses are bursty and
concentrated on few nodes, raising its rates to $2.0$--$31\%$.  With
identical evidence at all replicas, voting with two of five replicas
corrupted kept detection at the level of an honest twin,
whereas a corrupted single twin hid its node in every run.  The audit
log rejected every edit by an adversary holding fewer than a quorum of
replica keys.  Without congestion, an attacker that knows the rule can
lie in $37.5\%$ of rounds without being flagged.
\end{abstract}

\begin{keywords}
Audit log, digital twin, fog computing, Internet of Things, majority
voting, reputation, trust management.
\end{keywords}

\titlepgskip=-15pt

\maketitle

% ====================================================================
\section{Introduction}
\label{sec:intro}
% ====================================================================
\IEEEPARstart{T}{he} Internet of Things (IoT) is a network of
interconnected devices in which trust is fundamental to the security,
reliability and scalability of large deployments.  Individual devices
are cheap, constrained and routinely exposed to physical and network
compromise; a survey of IoT vulnerabilities, complemented by an
empirical measurement of exploited IoT devices on the Internet,
documents the practical consequences of this exposure~\cite{neshenko2019demystifying}.
Digital twin (DT) technology, reviewed for industry by Tao
\textit{et al.}~\cite{tao2019digital} and surveyed for the IoT by Minerva, Lee and
Crespi~\cite{minerva2020digital}, mirrors the state of each device so
that its activity can be cross-checked.

Soula \textit{et al.}~\cite{soula2024iodt} turned this idea into the
Internet of Digital Twins (IoDT) trust model.  IoDT assigns one
cloud-resident twin to every physical node and folds four binary
factors --- energy, packet delivery ratio (PDR), traceability and twin
notification --- into a single logical-AND verdict.  In its iFogSim
evaluation the model identifies the two compromised nodes in $92\%$
and $95\%$ of cases and isolates them within $250$--$300$~ms.  Three
limitations remain.

\begin{enumerate}
\item \emph{Fragility under transient congestion.}  Under the AND, a
single factor that drops for one round is enough to flag an honest
node.  The original evaluation already reports a false positive
(node N3) and notes that legitimate nodes can be mistaken for
compromised ones during collisions or high-traffic periods; its
authors suggest adjusting the threshold, or adapting the PDR
threshold under congestion, but do not evaluate such a variant~\cite{soula2024iodt}.
Dynamic-trust schemes~\cite{bao2012dynamic,nitti2014trustworthiness}
instead keep continuous scores that age out transient noise.

\item \emph{A single, corruptible twin and an untested adaptive
adversary.}  The original evaluation considers attackers that send
false notifications and a denial-of-service attacker~\cite{soula2024iodt}.
It does not consider on-off attackers that lie only occasionally to
stay below the detection threshold~\cite{sun2006trust},
nor an adversary that corrupts the twin that watches a node, in which
case nothing else in the pipeline can flag the node.  Threats against
the twin itself are part of the DT threat taxonomy
of~\cite{alcaraz2022digital}.

\item \emph{No tamper-evident record.}  Verdicts are stored in a
central database that can be silently rewritten if the cloud is
compromised.  Blockchain-based IoT and fog frameworks provide tamper
evidence~\cite{dorri2017blockchain,novo2018blockchain,kumar2020blockedge,%
dewanta2021bpt,akram2024tmiodt}, but most of them read from or write
to the chain while a decision is being made, and the latency and
computational overhead of blockchain validation is the reason Soula
\textit{et al.}~\cite{soula2024iodt} give for avoiding blockchain in
real-time IoT.
\end{enumerate}

This paper keeps the IoDT factor set and addresses the three
limitations with three extensions.  Its contributions are:
\begin{itemize}
\item a statement of IoDT trust evaluation as a per-round
classification problem, with metrics that separate per-round false
positives, the operationally relevant \emph{false-isolation} rate,
the mean time to the first false isolation over long operation, and
the cost of quarantine with re-admission (Section~\ref{sec:problem});
\item Extension~I, a weighted continuous score with exponential
moving-average (EMA) smoothing and a per-factor persistence check.
A closed form and a counting argument show that the binary AND errs at
rate $\Theta(q)$ in the per-factor congestion probability $q$, any
rule that needs two failures at $\Theta(q^{2})$, and the smoothed
score, at its default parameters and under independent congestion, at
$\Theta(q^{3})$.  The evaluation therefore compares it with
delay-matched baselines (a debounced AND and a 2-of-4 rule) and
quantifies the lie budget of a rule-aware on-off attacker;
\item Extension~II, quorum voting over signed, round-bound vote
vectors of replicated twins, under which, when the replicas share
their evidence and honest votes arrive in time, up to $k-1$ corrupted
replicas change no verdict, and missing or withheld votes never cause
a new isolation;
\item Extension~III, a signed Merkle audit log of verdicts, votes and
trust state that replicas endorse every $K$ rounds off the decision
path and that can be anchored in a permissioned ledger, with a
tamper-evidence theorem that makes explicit when an external
checkpoint is needed; and
\item a reproducible evaluation, with parameters selected on a
separate tuning seed, covering two calibrations to the original
evaluation, long horizons, stealthy, minimal-footprint and rule-aware
attackers, corrupted, withholding and slow replicas, bursty,
common-mode, cross-factor and uneven congestion, a $100$-sensor
topology, and a trace-driven evaluation on two public packet-delivery
datasets (Intel Lab and FlockLab).  The simulator, configurations, raw results and figure
scripts are provided with the submission.
\end{itemize}

The rest of the paper is organised as follows.
Section~\ref{sec:background} reviews related work,
Section~\ref{sec:problem} formalises the problem, the threat model and
the trust assumptions, Section~\ref{sec:proposedModel} presents the
three extensions with their analysis, Section~\ref{sec:resultsAndDiscussion}
reports the experiments, Section~\ref{sec:discussion} discusses them
and compares the design with related frameworks, and
Section~\ref{sec:conclusion} concludes.

% ====================================================================
\section{Related Work}
\label{sec:background}
% ====================================================================
\emph{Digital twins.}  Tao \textit{et al.}~\cite{tao2019digital}
review the state of the art of digital twins in industry, where a twin
integrates the physical and cyber spaces of a product or process.  Minerva \textit{et
al.}~\cite{minerva2020digital} widen the view to the IoT, and Sleuters
\textit{et al.}~\cite{sleuters2019digital} use twins for
the semi-automated behaviour analysis of large distributed IoT
systems.  Gehrmann
and Gunnarsson~\cite{gehrmann2020digital} propose a twin-based security
architecture for industrial control, Alcaraz and
Lopez~\cite{alcaraz2022digital} catalogue DT security threats in a
layered taxonomy that informs our threat model, and Empl \textit{et
al.}~\cite{empl2022soar4iot} prototype SOAR4IoT, a security orchestration,
automation and response application built on twins.

\emph{Digital twins and blockchain.}  Onwubiko \textit{et
al.}~\cite{onwubiko2023enabling} manage digital-twin documents with
Ethereum smart contracts and IPFS storage: stakeholders approve
document versions by a two-thirds threshold on chain and file hashes
are anchored for traceability; the prototype is evaluated by gas cost
and computes no trust score for devices.  Sasikumar \textit{et
al.}~\cite{sasikumar2023blockchain} connect twins of an industrial IoT
network through a proof-of-authority blockchain whose master nodes
validate the network, and report lower energy consumption and better
data security in simulation.  TMIoDT~\cite{akram2024tmiodt} combines a
digital-twin network, a blockchain-based credit score and a
deep-learning intrusion detector for drone services; it is discussed
in Section~\ref{sec:res-comparison}.

\emph{IoT trust management.}  Bao and Chen~\cite{bao2012dynamic} argue
for trust that adapts over time instead of a verdict computed once,
Nitti \textit{et al.}~\cite{nitti2014trustworthiness} carry the idea
to the social IoT, and EigenTrust~\cite{kamvar2003eigentrust} gives the
canonical recipe for continuous reputation by power iteration.  The
beta reputation system of J{\o}sang and Ismail~\cite{josang2002beta}
updates trust from counts of positive and negative observations with
a forgetting factor.  Sun \textit{et
al.}~\cite{sun2006trust} analyse on-off and bad-mouthing attacks on
trust evaluation, and Chen, Bao and Guo~\cite{chen2016trust} evaluate
trust management for the social IoT under self-promoting,
bad-mouthing, ballot-stuffing and discriminatory attacks.  The
surveys~\cite{yan2014survey,guo2017survey,konsta2023survey} review the
field.  Hussain \textit{et al.}~\cite{hussain2020context} propose a
context-aware trust and reputation model for fog-based IoT in which a
device's context is its role or purpose in the network and the trust
level in $[0,1]$ is updated recursively from previous values, with a
weight of $0.7$ on the current trust level, and multi-source feedback.

\emph{Fog and edge computing.}  Bonomi \textit{et
al.}~\cite{bonomi2012fog} introduce fog computing, Chiang and
Zhang~\cite{chiang2016fog} survey its research opportunities for the
IoT, and Shi \textit{et
al.}~\cite{shi2016edge} push computation further to the edge.

\emph{Blockchain for IoT and fog.}  Dorri \textit{et
al.}~\cite{dorri2017blockchain} build a lightweight chain for the smart
home, and Novo~\cite{novo2018blockchain} uses one for scalable access
management.  Debe \textit{et
al.}~\cite{debe2019iot} add an Ethereum-based reputation system for
public fog nodes.  BlockEdge~\cite{kumar2020blockedge} integrates
blockchain pre-processing into an industrial edge pipeline, and
BPT~\cite{dewanta2021bpt} establishes trust between vehicles before a
fog-computing service through a bid and rating exchange recorded on a
chain.

\emph{Tamper-evident logs and anchoring.}  Tamper evidence does not
require a blockchain.  Haber and Stornetta~\cite{haber1991time} link
time-stamped documents by hashes, Schneier and
Kelsey~\cite{schneier1999secure} make audit logs on untrusted machines
tamper-evident with hash chains, and Crosby and
Wallach~\cite{crosby2009efficient} use history trees to make such logs
efficient to audit.  Certificate
Transparency~\cite{laurie2013rfc6962} publishes signed Merkle tree
heads whose leaf and interior hashes are domain-separated.  The
Keyless Signature Infrastructure of Buldas \textit{et
al.}~\cite{buldas2013ksi} publishes only the root of a hash tree, and
Arbitrum~\cite{kalodner2018arbitrum} executes smart contracts
off-chain while the chain records only hashes of their state.
Extension~III follows this line: it is a signed Merkle log that a
ledger may host but whose guarantee does not depend on one.
Section~\ref{sec:res-comparison} compares the closest trust systems on
the figures they actually report.

\emph{Network-level anomaly detection.}  N-BaIoT~\cite{meidan2018nbaiot}
detects botnet traffic with deep autoencoders.  It observes traffic
rather than behavioural trust factors and is complementary to the
present work.

% ====================================================================
\section{Problem Statement}
\label{sec:problem}
% ====================================================================
For every pair of communicating nodes, the IoDT model
of~\cite{soula2024iodt} evaluates four binary factors.  Energy $E$ is
one when the energy a node consumed lies within the bounds
$[E_{\min},E_{\max}]$; $E_{\max}$ follows from the capacity of the
node's super-capacitor and $E_{\min}$ from the CSMA/CA transmission
and back-off cost, both scaled by the number of packets sent.  The PDR
score is one when the measured PDR reaches a threshold
$\Theta_{\mathrm{PDR}}$ ($90\%$ is the example value
of~\cite{soula2024iodt}).  Traceability $T$ is one when the node's most
recent claimed transaction is present in the trust database, and the
twin-notification score $W$ is one unless the twin of the destination
named in the node's notification denounces the claimed communication
(our reading of the notification check of~\cite{soula2024iodt}).  The
verdict is the logical AND
\begin{equation}
\mathrm{Trust}(i) =
E(i) \wedge \mathrm{PDR}(i) \wedge T(i) \wedge W(i).
\label{eq:binary_and}
\end{equation}

\subsection{System Notation}
\label{subsec:notation}
Let $\mathcal{N}$ be the set of physical nodes and
$\mathcal{N}_c\subset\mathcal{N}$ the unknown set of compromised
nodes.  Time is divided into evaluation rounds $t=1,2,\dots$; in every
round the trust engine computes, for every node $i$, the factor vector
\begin{equation}
\mathbf{x}_t(i) =
\bigl(E_t,P_t,T_t,W_t\bigr)(i)
\in \{0,1\}\times[0,100]\times\{0,1\}^{2}
\label{eq:factor_vector}
\end{equation}
where $P_t$ is the PDR in percent.  A trust-aggregation function $f$
maps the history $\mathbf{x}_{1:t}(i)$ to a verdict
$\hat{y}_t(i)\in\{0,1\}$, with $\hat{y}=0$ meaning ``flagged''.

\subsection{Performance Metrics}
\label{subsec:metrics}
The following metrics are used throughout.
\begin{itemize}
\item The \emph{per-round false-positive rate} over a horizon of $T$
rounds,
\begin{equation}
\mathrm{FPR}=\frac{1}{T\,|\mathcal{N}\setminus\mathcal{N}_c|}
\sum_{i\notin\mathcal{N}_c}\sum_{t=1}^{T}\mathbb{1}[\hat{y}_t(i)=0].
\label{eq:fpr}
\end{equation}
\item The \emph{false-isolation rate}: the share of honest nodes
flagged at least once within the horizon $T$.  IoDT isolates a node as
soon as it is flagged~\cite{soula2024iodt}, so this is the number that
matters for a run.  It grows with $T$; the main tables use $T=16$
rounds, the length of one reference scenario
(Section~\ref{sec:setup}), and Section~\ref{sec:res-long} reports it
for up to $1{,}024$ rounds.
\item The \emph{mean time to first false isolation} (MTTFI) of an
honest node over long operation, estimated as the inverse of the rate
at which flags start on honest nodes; it does not depend on $T$.
Strictly, this is the mean time between flag onsets; when flags are
frequent it slightly overstates the time to the first one (about $8$
against $7$ rounds for the AND at $q=0.04$), which does not change any
comparison below.
\item Under \emph{quarantine with re-admission}, a flagged node is
quarantined and re-admitted after $R$ consecutive rounds in which a
quorum of its twins reports $\tau_t\ge\theta_{r}$ and no persistence
trigger ($R=8$, $\theta_r=0.9$); for the baselines, which have no
$\tau_t$, the condition is $R$ consecutive rounds without a flag.  We report the share of time an honest
node spends in quarantine, the number of false quarantines per node
and hour, and the share of rounds in which an attacker delivers a lie
while it is not quarantined.
\item The \emph{detection rate}: the share of runs in which a
compromised node is flagged at least once within $T$ rounds, together
with the mean \emph{time to detection} (TTD) in rounds, and the
\emph{steady-state per-round detection}, the share of rounds $t\geq3$
in which a compromised node is flagged.
\end{itemize}
\subsection{Threat Model}
\label{subsec:threat}
Following the IoT threat survey of~\cite{neshenko2019demystifying} and
the DT threat taxonomy of~\cite{alcaraz2022digital}, the adversary can
compromise a subset $\mathcal{N}_c$ of physical nodes, corrupt at most
$c$ of the $n$ twin replicas that watch a node, steal replica signing
keys, and rewrite or delete records in the cloud.  Compromised nodes
behave in one of four ways.
\begin{itemize}
\item \emph{Lie (IoDT scenario).}  As in~\cite{soula2024iodt}, the
attacker reports a false destination to its twin, and the twin of the
falsely named destination denounces the claim.  We model a lie as
setting $W=0$ and, because the claimed transaction is not found in the
database, also $T=0$; the second effect is our modelling choice.  The
attacker lies in each round with probability $p\in(0,1]$; $p=1$ is the
original naive attacker and $p<1$ a stealthy one.
\item \emph{Minimal-footprint lie.}  The attacker arranges for the
claimed transaction to be logged, so a lie costs only $W=0$, the
smallest observable deviation.
\item \emph{Rule-aware on-off lie.}  The attacker knows the
aggregation rule and its parameters, observes its own past evidence,
and lies only in rounds in which, absent congestion, no trigger would
fire; once flagged it behaves honestly until it is re-admitted.  This
is the on-off attack of~\cite{sun2006trust} against a known detector.
\item \emph{Corrupted twins.}  The adversary controls $c$ replicas
of the twins that watch the two compromised nodes.  Corrupted replicas
vote \textsc{trust} for the attackers (\emph{shielding}), may also
vote \textsc{flag} for honest nodes (\emph{framing}), or may withhold
their votes.
\end{itemize}
Attackers that never deviate on any factor are indistinguishable from
honest nodes at this layer and are out of scope.

Table~\ref{tab:threats} lists which component the design trusts for
what, and which attacks from the IoT trust literature it covers.

\begin{table}[!t]
\caption{Trust assumptions and threat coverage.}
\label{tab:threats}
\centering
\renewcommand{\arraystretch}{1.1}
\footnotesize
\setlength{\tabcolsep}{2.5pt}
\begin{tabular}{p{0.30\columnwidth} p{0.62\columnwidth}}
\toprule
\textbf{Component / attack} & \textbf{Assumption or coverage} \\
\midrule
Fog relay & Trusted to relay evidence and to enforce isolation; a compromised fog node is out of scope \\
Twin replicas & Up to $k-1$ of $n$ Byzantine: shield, frame or withhold (Sec.~\ref{sec:res-ext2}) \\
Trust engine & Trusted for timeliness; the enforced verdicts are checked against the signed votes at every header commit, so an unsupported verdict is exposed within $K$ rounds plus the append time, but not prevented (Sec.~\ref{sec:ext2}) \\
Log storage, ledger & May delete or rewrite; rewrites are rejected by the verifier (Theorem~\ref{thm:tamper}); deleting the newest windows is detected only against a newer checkpoint \\
Checkpoint channel & Controlled independently of the cloud operator \\
\midrule
Lies, stealthy lies & Covered for frequent lies; liars at $p\le0.25$ are detected slowly (Sec.~\ref{sec:res-ext1}, \ref{sec:res-stealth}) \\
Rule-aware on-off & Partly: a lie budget remains (Sec.~\ref{sec:res-stealth}) \\
Bad-mouthing via $W$ & Not covered: a compromised destination can denounce an honest neighbour; repeated denunciations trigger the persistence check \\
Ballot-stuffing via $W$ & Partly: a colluding destination can hide the $W$ failure; under our lie model the lie still fails $T$, but a liar that also gets its transaction logged is not detected \\
Sybil, whitewashing & Out of scope: identities are provisioned by the operator~\cite{douceur2002sybil,friedman2001social} \\
Denial of service & Detected in~\cite{soula2024iodt} mainly through abnormal energy use (a DoS against a twin); not re-evaluated here \\
\bottomrule
\end{tabular}
\end{table}

% ====================================================================
\section{The Proposed Trust Model}
\label{sec:proposedModel}
% ====================================================================
This section presents the three extensions to the IoDT
model~\cite{soula2024iodt}, whose architecture is shown in
Fig.~\ref{fig:architecture}.  Twin replicas run Extension~I, a quorum
over their signed votes (Extension~II) decides on the real-time path,
and verdicts and votes are committed to a signed Merkle audit log off
that path (Extension~III).

\Figure[!t](topskip=0pt,botskip=0pt,midskip=0pt)[width=\textwidth]{figs/fig_architecture.png}
{Architecture of the proposed IoDT trust pipeline with its three extensions.\label{fig:architecture}}

\subsection{Extension I: Continuous Trust Score}
\label{subsec:ext1-math}
The four factors are mapped to $[0,1]$ and combined with non-negative
weights that sum to one:
\begin{equation}
s_{t}(i) = w_{E}E_{t}(i) + w_{P}\frac{P_{t}(i)}{100}
         + w_{T}T_{t}(i) + w_{W}W_{t}(i).
\label{eq:weighted_score}
\end{equation}
All experiments use $w_{E}=w_{P}=0.30$ and $w_{T}=w_{W}=0.20$, which
put slightly more weight on the physical factors than on the
notification-driven ones.  A single bad round should not condemn a
node that has otherwise behaved well, so the score is smoothed by an
exponential moving average,
\begin{equation}
\tau_{t}(i) = \alpha\,s_{t}(i) + (1-\alpha)\,\tau_{t-1}(i),
\qquad \tau_{0}(i)=1,
\label{eq:ema}
\end{equation}
and the node is flagged when $\tau_t(i)<\theta$.  The defaults are
$\alpha=0.5$ and $\theta=0.70$, selected on a separate tuning seed
(Section~\ref{sec:setup}).

Smoothing has a blind spot: a weighted score cannot react to an
attacker that always fails the \emph{same} low-weight factor.  An
attacker that fails only $W$ keeps $s_t\in[0.776,0.80]$, and, absent
congestion, its EMA never falls below $0.70$.  Extension~I therefore adds a
\emph{per-factor persistence check}: the node is also flagged when any
single factor has failed in at least $m$ of the last $w$ rounds.
Congestion scatters failures across factors and rounds, whereas a
liar fails the same factor again and again.  The defaults are $m=4$
and $w=8$.  A second, longer window ($m_2$ of $w_2$ rounds) can be
added against rule-aware attackers (Section~\ref{sec:res-stealth}); it
is off by default.  Algorithm~\ref{alg:ema} gives the complete
per-round update.

\begin{algorithm}[!t]
\caption{Extension~I trust update for node $i$ at round $t$.}
\label{alg:ema}
\noindent\textbf{Input:} factor scores $(E_t,P_t,T_t,W_t)$; weights;
prior trust $\tau_{t-1}(i)$; failure counters of the last $w$ (and
$w_2$) rounds; $\alpha$, $\theta$, $m$, $w$ (and $m_2$, $w_2$).

\noindent\textbf{Procedure:}
\begin{enumerate}\setlength{\itemsep}{0pt}\setlength{\parskip}{0pt}
  \item $s_t \gets w_E E_t + w_P (P_t/100) + w_T T_t + w_W W_t$.
  \item $\tau_t(i) \gets \alpha\,s_t + (1-\alpha)\,\tau_{t-1}(i)$.
  \item Update, for each factor, the number of failures in rounds
        $t-w+1,\dots,t$ (a factor fails when it is $0$, or when
        $P_t<\Theta_{\mathrm{PDR}}$).
  \item If $\tau_t(i)<\theta$, or some factor failed in $\ge m$ of the
        last $w$ rounds (or in $\ge m_2$ of the last $w_2$), return
        \textsc{flag}; else return \textsc{trust}.
\end{enumerate}
\end{algorithm}

\subsubsection{False-positive analysis}
The analysis rests on three assumptions, and the experiments test what
happens when the first two are relaxed.
\begin{itemize}
\item[(A1)] \emph{Independence.}  For an honest node, factor
corruptions are independent across factors and rounds.
\item[(A2)] \emph{Symmetric corruption.}  Each factor is corrupted
with probability $q$; a corrupted binary factor becomes $0$ and a
corrupted PDR falls below $\Theta_{\mathrm{PDR}}$.
\item[(A3)] \emph{Healthy-link PDR.}  An uncorrupted PDR lies in
$[\underline{P},100]$; the simulator uses $\underline{P}=92\%$.
\end{itemize}
Under (A1)--(A2) the binary AND (\ref{eq:binary_and}) flags an honest
node whenever any factor is corrupted, so its per-round FPR is
\begin{equation}
\mathrm{FPR}_{\mathrm{AND}}(q) = 1-(1-q)^{4} = 4q + O(q^{2}),
\label{eq:fpr_and}
\end{equation}
and by the same counting any rule that flags only after two failures
--- the AND required to fail in two consecutive rounds, or two of the
four factors required to fail in the same round --- errs at rate
$\Theta(q^{2})$.  In the evaluation the 2-of-4 rule is combined with
the persistence check below, which keeps its order at $\Theta(q^{2})$
and lets it see a liar that fails only one factor.  These rules serve as delay-matched baselines in
Section~\ref{sec:resultsAndDiscussion}.  The EMA needs three failures.

\begin{proposition}[FPR of the EMA score]\label{prop:fpr_ema}
A node is flagged by the EMA at round $t$ if and only if
\begin{equation}
s_t < s^{\star}=\frac{\theta-(1-\alpha)\,\tau_{t-1}}{\alpha}.
\label{eq:sstar}
\end{equation}
With the default weights, $\alpha=0.5$ and $\theta=0.70$, a flag
requires at least three factor failures within the last $L=6$ rounds,
so under (A1)--(A3)
\begin{equation}
\mathrm{FPR}_{\mathrm{EMA}}(q) \le \binom{4L}{3}q^{3},\qquad
\mathrm{FPR}_{\mathrm{EMA}}(q)=\Theta(q^{3}).
\end{equation}
\end{proposition}
\begin{proof}
See Appendix~\ref{app:proofs}.
\end{proof}

The cubic order depends on the parameters and on (A1).  At
$\alpha\geq0.6$ (with $\theta=0.70$) or $\theta\geq0.75$ (with
$\alpha=0.5$), two failures in one round can already cross the
threshold and the order falls to $\Theta(q^{2})$
(Appendix~\ref{app:proofs}).  When one congestion event fails several
factors of a node at once, (A1) fails across factors, and the order is
again $\Theta(q^{2})$ (Section~\ref{sec:res-robust}).  The price of
smoothing is one round of delay: without congestion, a node that lies
in every round scores $s_t<0.60$, so from $\tau_0=1$ its EMA is below
$0.70$ after two rounds, whereas the AND flags it in the first.

The persistence check adds false positives at a controlled rate and
catches lies that the EMA cannot see.

\begin{proposition}[Persistence check]\label{prop:persist}
Under (A1)--(A2) the probability that the persistence check flags an
honest node in a given round is at most
$4\binom{w}{m}q^{m}$ ($7.2\times10^{-4}$ at $m=4$, $w=8$, $q=0.04$).
A node whose lies fail a given factor independently in each round with
probability $p$ is flagged within any window of $w$ rounds with
probability at least $\Pr[\mathrm{Bin}(w,p)\geq m]$; for $p=1$ it is
flagged by round $m$.  A node that keeps at most $m-1$ lies in every
window of $w$ rounds is never flagged by the check: its \emph{lie
budget} is $(m-1)/w$, or $\min\{(m-1)/w,(m_2-1)/w_2\}$ with a second
window.
\end{proposition}
\begin{proof}
See Appendix~\ref{app:proofs}.
\end{proof}

% ====================================================================
\subsection{Extension II: Multi-Twin Voting}
\label{sec:ext2}
% ====================================================================
A single twin is a single point of trust: if it is corrupted, the
node it watches is never flagged.  Extension~II deploys $n$ replicas
of each twin and decides by a quorum of their signed votes; the
reported configuration is $n=5$, $k=3$.  This is one-shot voting over
signed verdicts in the sense of Byzantine quorum
systems~\cite{malkhi1998byzantine}, not a full Byzantine-fault-tolerant
agreement protocol such as PBFT~\cite{castro2002practical}: there is no
multi-phase commit, view change or total order on operations.  The
extension relies on the following assumptions.
\begin{itemize}
\item[(B1)] \emph{Independent hosting.}  The replicas run on
administratively independent hosts (different fog sites or cloud
availability zones), so compromising one host does not compromise the
others.
\item[(B2)] \emph{Authenticated, bound votes.}  In every round each
replica signs one vector of verdicts for all the nodes it watches,
bound to its identity, the round number and the current key epoch
(ECDSA over secp256k1 in our implementation).  Keys are rotated by
starting a new epoch, which invalidates signatures from retired keys.
\item[(B3)] \emph{Bounded delay.}  Votes from honest replicas arrive
within a known bound $\Delta$; a vote that arrives later is treated as
missing.
\item[(B4)] \emph{Identical evidence.}  The fog node forwards the same
evidence to every replica, so honest replicas reach identical
verdicts.
\item[(B5)] \emph{Adversary budget.}  At most $c\le k-1$ replicas are
corrupted.
\end{itemize}
Each replica runs Algorithm~\ref{alg:ema}, signs its vote vector and
sends it to the trust engine, which verifies the signatures and
applies the \emph{quorum rule} of Algorithm~\ref{alg:vote}:
\textsc{trust} if at least $k$ valid votes say \textsc{trust},
\textsc{flag} if at least $n-k+1$ say \textsc{flag}, and otherwise the
previous verdict is kept and an availability alarm raised.
The engine forwards the verdict to the fog node, which enforces it
and, being trusted (Table~\ref{tab:threats}), appends the verdict it
enforced to the off-chain store of Extension~III.  With all
$n=2k-1$ votes present exactly one quorum is reached, so the rule
coincides with ``flag if fewer than $k$ \textsc{trust} votes''; unlike
that rule, missing or withheld votes never cause a new isolation,
although a node that is already isolated stays isolated while no
quorum forms.  The count is computed by the trust engine, not by a smart
contract, so no ledger operation lies on the decision path.  The
verdict and the signed votes are appended to the off-chain store of
Extension~III.

\begin{algorithm}[!t]
\caption{Per-round vote for node $i$.}
\label{alg:vote}
\noindent\textbf{Input:} evidence $\mathbf{x}_t(i)$; replicas
$\mathrm{DT}_1,\dots,\mathrm{DT}_n$; quorum $k$; previous verdict.

\noindent\textbf{Procedure:}
\begin{enumerate}\setlength{\itemsep}{0pt}\setlength{\parskip}{0pt}
  \item Each $\mathrm{DT}_r$ runs Algorithm~\ref{alg:ema} on
        $\mathbf{x}_t(i)$ and signs its vote vector ($\sigma_r$),
        bound to $(r,t,\text{epoch})$.
  \item The trust engine waits at most $\Delta$, verifies each received
        signature and its binding, and counts the \textsc{trust} votes
        $V_T$ and the \textsc{flag} votes $V_F$.
  \item If $V_T\geq k$ return \textsc{trust}; if $V_F\geq n-k+1$ return
        \textsc{flag} to the fog node, which isolates node $i$;
        otherwise keep the previous verdict and raise an alarm.
  \item Append the verdict and the signed votes to the off-chain store
        (input to Algorithm~\ref{alg:anchor}).
\end{enumerate}
\end{algorithm}

Under (B3)--(B5) the quorum behaves like one honest twin.  Corrupted
replicas that always vote \textsc{trust} for an attacker
(\emph{shielding}) are outvoted whenever the $n-c\geq n-k+1$ honest
replicas flag it, and corrupted replicas that vote \textsc{flag} for an
honest node (\emph{framing}) are outvoted whenever the honest replicas
trust it.  Detection and false positives are therefore those of an
honest single twin for every $c\le k-1$, whereas a single corrupted
twin never flags its node and $c\geq k$ corrupted replicas decide every
verdict.

The trust engine itself is not replicated.  It is trusted to count
promptly, but not trusted for integrity: every enforced verdict is
logged with the signed votes behind it, and replicas endorse a header
of the audit log only if each verdict in it follows from its votes by
the quorum rule (Section~\ref{sec:ext3}).  An engine that issues an
unsupported verdict is therefore exposed when the next header is
committed, within $K$ rounds plus the time to append it, although it
is not prevented in real time.  Voting defends
against corrupted twins; it does not help against attackers that lie
rarely, whose detection is governed by
Eq.~(\ref{eq:sstar}) and Proposition~\ref{prop:persist}.
If the replicas receive independent evidence instead of (B4), framing
replicas can tip honest nodes whose honest replicas disagree
(Section~\ref{sec:res-ext2}).

% ====================================================================
\subsection{Extension III: Signed Merkle Audit Log}
\label{sec:ext3}
% ====================================================================
Most blockchain-based trust schemes consult or update the chain while
a decision is being made (Section~\ref{sec:res-comparison}).  Here the
verdict is produced by Extensions~I and~II alone, and an audit log
records it afterwards.  The log is a hash-chained sequence of signed
headers over Merkle-committed off-chain data, in the tradition of
tamper-evident logs~\cite{schneier1999secure,crosby2009efficient,laurie2013rfc6962}.
Its headers can be kept in any append-only store; hosting them on a
permissioned ledger (for example Hyperledger Fabric with a
crash-tolerant ordering service among fog peers) adds replicated
storage and availability against an operator that deletes the log, but
the tamper evidence of Theorem~\ref{thm:tamper} does not depend on the
ledger's consensus.

\emph{Content.}  Every round, the per-node event (factor scores,
$\tau_t$, verdict; $|E|=29$~bytes) and the replicas' signed vote
vectors are stored off-chain.  Every $K$ rounds the trust engine
proposes a header that carries the height, the first and last round,
the numbers of events and votes, the key epoch, the previous header
hash, and the Merkle roots of the window's events, of its signed votes
and of the per-node state records ($|R|=13$~bytes per node);
$|H|=152$~bytes.  The trees follow RFC~6962~\cite{laurie2013rfc6962}:
leaf and interior hashes are domain-separated, the last node is not
duplicated, and the leaf count is committed in the header.  Each
replica \emph{endorses} the header --- signs it --- only after checking
the off-chain data against it, including that every verdict follows
from its votes; a header needs $k$ endorsements.  A partial final
window is committed as well.  Each header therefore occupies
$|H|+k\,|\sigma|$ bytes in the header store, with $|\sigma|\le73$~bytes per
endorsement (identifier and DER signature), independent of the number
of nodes.  Raw evidence never
leaves the off-chain store, which also limits what the log reveals
about each node.  A verdict becomes auditable within $K$ rounds plus
the time to append a header; this delay never affects when the verdict
itself is issued.  We further assume:
\begin{itemize}
\item[(C1)] header hashes are periodically published on an
independently controlled channel (linked
time-stamping~\cite{haber1991time}, a KSI hash
tree~\cite{buldas2013ksi}, an Arbitrum-style state
commitment~\cite{kalodner2018arbitrum}, or a public-chain anchor);
\item[(C2)] the signature scheme is existentially unforgeable and
the hash function is collision resistant.
\end{itemize}

\begin{algorithm}[!t]
\caption{Audit-log commit at round $t$.}
\label{alg:anchor}
\noindent\textbf{Input:} pending events $\mathcal{E}$ and votes
$\mathcal{S}$; state records; period $K$; previous header hash
$H_{\mathrm{prev}}$.

\noindent\textbf{Procedure:}
\begin{enumerate}\setlength{\itemsep}{0pt}\setlength{\parskip}{0pt}
  \item Append the round's events and signed votes (off-chain).
  \item If $t \bmod K \neq 0$ and the run is not ending, return.
  \item The trust engine forms the header with
        $\mathrm{MTH}(\mathcal{E})$, $\mathrm{MTH}(\mathcal{S})$,
        $\mathrm{MTH}(\text{state})$, counts, epoch and
        $H_{\mathrm{prev}}$.
  \item Each replica checks the window (lengths, counts, roots, vote
        signatures and bindings, verdicts against votes, state against
        events) and, if it passes, signs the header.
  \item With $k$ endorsements, append the header; periodically
        publish its hash externally (C1).
\end{enumerate}
\end{algorithm}

\begin{definition}[Verifier]
Given the headers and the off-chain data, the verifier $V$ checks the
hash links, at least $k$ valid endorsements from distinct replicas of
the header's key epoch, record lengths and counts, every Merkle root,
every vote signature and its binding to replica, round and epoch, that
each verdict follows from its round's votes by the quorum rule, that
each state record equals the node's latest event, and, if a checkpoint
$(j,h^{\star}_j)$ is available, that header $j$ hashes to $h^{\star}_j$.
\end{definition}

\begin{theorem}[Tamper evidence]\label{thm:tamper}
Under (C2), let the adversary modify, insert, delete or reorder any
event, vote or state record of a committed window $B_i$ and recompute
any hashes it likes.
\begin{enumerate}
\item If it holds fewer than $k$ replica signing keys, $V$ rejects
except with negligible probability.
\item If it holds $k$ or more keys, it can make $V$ accept if and only
if no checkpoint $(j,h^{\star}_j)$ with $j\ge i$ was published before
the modification (up to a hash collision).
\end{enumerate}
\end{theorem}
\begin{proof}
See Appendix~\ref{app:proofs}.
\end{proof}

Endorsements, not the checkpoint, are therefore the first line of
defence; the checkpoint is what defeats an adversary that steals or
colludes with $k$ replicas.

\subsection{End-to-End Round}
\label{sec:round}
Fig.~\ref{fig:sequence} shows how the three extensions interact in one
round.  Steps 1--7 form the real-time decision path: the fog node
forwards the evidence to the replicas, each replica runs
Algorithm~\ref{alg:ema} and signs its vote vector, and the trust engine
applies Algorithm~\ref{alg:vote} and returns the verdict to the fog
node, which isolates a flagged node.  Steps 8--11 run afterwards and
never delay a verdict: events and votes are appended to the off-chain
store, and every $K$ rounds Algorithm~\ref{alg:anchor} forms a header,
collects $k$ replica endorsements and appends it; header hashes are
periodically published on the checkpoint channel (C1).

\Figure[!t](topskip=0pt,botskip=0pt,midskip=0pt)[width=0.92\textwidth]{figs/fig_sequence.png}
{Message sequence of one round: the decision path (steps 1--7) and the off-path audit-log commit (steps 8--11).\label{fig:sequence}}

% ====================================================================
\section{Results}
\label{sec:resultsAndDiscussion}
% ====================================================================

\subsection{Experimental Setup}
\label{sec:setup}
The extensions are evaluated with a discrete-round Python simulator of
the IoDT factor model.  It follows the topology of the reference
scenario of~\cite{soula2024iodt}: eight sensors, two of which
(COMP1 and COMP2, labelled C1 and C2 in Fig.~\ref{fig:architecture}) are compromised, attached to two fog nodes, with
twins and the trust engine in the cloud.  The simulator models the
trust logic only; it is not a re-implementation of the iFogSim
evaluation of~\cite{soula2024iodt}, whose resource and energy model it
does not include~\cite{gupta2017ifogsim}.  Delays are counted in
evaluation rounds.  Where a time is given, it assumes a round length
$T_{\mathrm{round}}=117$~ms, obtained by dividing the $1875$-ms
execution time reported for the IoDT scenario in~\cite{soula2024iodt}
into 16 rounds; the
original model itself has no rounds, so this conversion is an
assumption.  A run lasts $T=16$ rounds unless stated otherwise.

In every round each factor of every node is corrupted independently
with probability $q$ (assumption A1).  Because the congestion level of
a real deployment is unknown, results are reported over a range of
$q$.  Two levels are tied to the original evaluation, in which one of
six honest nodes (N3) was misidentified in $2\%$ of
cases~\cite{soula2024iodt}, a figure whose unit is not stated.
\emph{Reading~A} takes it as a per-node, per-round rate of $2\%/6$,
which the binary AND approximately reaches at $q=0.001$ (per-round
FPR $0.40\%$ against $0.33\%$).  \emph{Reading~B} takes it as a
per-run false-isolation rate over the 16-round scenario, again
divided among the six honest nodes ($2\%/6$ per node), which the AND
reaches at $q=5.2\times10^{-5}$.  Neither calibration
reproduces the detection side: the simulated AND detects both attackers
in every run, against $92\%$ and $95\%$ in~\cite{soula2024iodt},
because a lie here always fails $W$ and $T$, while the causes of the
missed detections in the iFogSim runs are not reported.  The stress
level is $q=0.04$.  A healthy PDR is drawn uniformly from $[92,100]\%$
and a degraded one from $[60,89]\%$, against the $90\%$ threshold.

Extension~I is compared with the binary AND of~\cite{soula2024iodt}
and with two delay-matched rules: the AND required to fail in two
consecutive rounds (\emph{debounced AND}), which has the same one-round
delay as Extension~I, and a rule that flags when two of the four
factors fail in the same round or the persistence check fires
(\emph{2-of-4 rule}).  A beta reputation score~\cite{josang2002beta}
on the AND observation, tuned on seed~1 by an analogous rule
(naive liar flagged within two rounds without congestion, lowest FPR
at $q=0.04$, near-ties resolved towards the simpler setting),
reduced to the debounced AND (its FPR was within $0.002$ percentage
points in every cell of Table~\ref{tab:headline}) and is not shown.
The \emph{full system} is Extension~I in each of five replicas with a
$3$-of-$5$ quorum.  Because its replicas receive identical evidence
(B4), its verdicts coincided with those of Extension~I in every cell
of Table~\ref{tab:headline} (its row is therefore omitted), also with two corrupted replicas that
shield the attackers and frame honest nodes;
Section~\ref{sec:res-ext2} evaluates corrupted replicas in detail.

The parameters of Extension~I were selected on a tuning seed (seed~$1$)
by two rules: the lowest per-round FPR at $q=0.04$ among settings that
flag a naive liar within two rounds on average, which selects
$(\alpha,\theta)=(0.5,0.70)$; and the lowest FPR among persistence
windows that flag a persistent minimal-footprint liar within four
rounds and detect one that lies in half of the rounds in at least
$90\%$ of runs, which selects $4$ of $8$.  Every reported cell uses
seed~$0$ with $2{,}000$ runs unless stated otherwise (the long-horizon
runs use seeds $100$--$105$).  Per-round rates carry $95\%$
percentile-bootstrap intervals over runs, which respect the correlation
between rounds of a run, and run-level proportions carry $95\%$ Wilson
intervals.  For a zero count we give the one-sided $95\%$ bound $3/N$
over the $N$ honest node-runs ($0.025\%$ for $N=12{,}000$), which also
bounds the per-round FPR.  Table~\ref{tab:params} lists the defaults.
The experiments ran in Python~3.13 with NumPy and the
\texttt{cryptography} package (OpenSSL) on an Intel Xeon at $2.8$~GHz,
single-threaded.  The tuning grids are provided with the
supplementary material.

\begin{table}[!t]
\caption{Default parameters of the evaluation.}
\label{tab:params}
\centering
\renewcommand{\arraystretch}{1.1}
\footnotesize
\setlength{\tabcolsep}{3pt}
\begin{tabular}{lll}
\toprule
\textbf{Symbol} & \textbf{Meaning} & \textbf{Default} \\
\midrule
$w_E,w_P,w_T,w_W$ & weights (\ref{eq:weighted_score}) & $.3,.3,.2,.2$ \\
$\alpha,\ \theta$ & EMA factor, threshold (\ref{eq:ema}) & $0.5,\ 0.70$ \\
$m,\ w$ & persistence check & $4$ of $8$ rounds \\
$m_2,\ w_2$ & optional second window & off ($8$ of $32$) \\
$q$ & per-factor congestion rate & $5.2\times10^{-5}$--$0.04$ \\
$T$ & horizon of a run & $16$ rounds \\
$R,\ \theta_r$ & re-admission after & $8$ rounds with $\tau_t\ge0.9$ \\
$n,\ k$ & replicas, quorum & $5,\ 3$ \\
$K$ & header (commit) period & $10$ rounds \\
Runs & Monte-Carlo runs per cell & $2{,}000$ \\
\bottomrule
\end{tabular}
\end{table}

% ====================================================================
\subsection{Congestion Tolerance}
\label{sec:res-ext1}
Table~\ref{tab:headline} compares the aggregators at five congestion
levels.  Every aggregator detected both naive attackers in every run at
every level, so the table reports what differs: false positives and
delay.

\begin{table*}[!t]
\caption{Congestion tolerance of Extension~I and the baselines (2{,}000 runs, 16 rounds, two significant figures; 95\% intervals in brackets; zero counts with their one-sided 95\% bound).}
\label{tab:headline}
\centering
\renewcommand{\arraystretch}{1.05}
\setlength{\tabcolsep}{5pt}
\footnotesize
\begin{tabular}{lrrrrr}
\toprule
 & & & \textbf{TTD, naive} & \multicolumn{2}{c}{\textbf{Minimal-footprint liar}} \\
\cmidrule(lr){5-6}
\textbf{Aggregator} & \textbf{Per-round FPR \%} & \textbf{False isolation \%} &
\textbf{liar (rounds)} & \textbf{Det.\ \%} & \textbf{TTD (rounds)} \\
\midrule
\multicolumn{6}{l}{$q=5.2\times10^{-5}$ (reading B of~\cite{soula2024iodt})} \\
Binary AND~\cite{soula2024iodt} & 0.024 [0.018,\,0.032] & 0.39 [0.29,\,0.52] & 1.00 & 100 & 1.0 \\
Debounced AND & 0 ($<$0.025) & 0 ($<$0.025) & 2.00 & 100 & 2.0 \\
2-of-4 rule & 0 ($<$0.025) & 0 ($<$0.025) & 1.00 & 100 & 4.0 \\
\textbf{Extension~I} & 0 ($<$0.025) & 0 ($<$0.025) & 2.00 & 100 & 4.0 \\
\midrule
\multicolumn{6}{l}{$q=0.001$ (reading A of~\cite{soula2024iodt})} \\
Binary AND~\cite{soula2024iodt} & 0.40 [0.37,\,0.43] & 6.2 [5.8,\,6.6] & 1.00 & 100 & 1.0 \\
Debounced AND & 0.0010 [0,\,0.0026] & 0.017 [0.0046,\,0.061] & 2.00 & 100 & 2.0 \\
2-of-4 rule & 0.00052 [0,\,0.0016] & 0.0083 [0.0015,\,0.047] & 1.00 & 100 & 4.0 \\
\textbf{Extension~I} & 0 ($<$0.025) & 0 ($<$0.025) & 2.00 & 100 & 4.0 \\
\midrule
\multicolumn{6}{l}{$q=0.005$} \\
Binary AND~\cite{soula2024iodt} & 2.0 [1.9,\,2.0] & 27 [26,\,28] & 1.00 & 100 & 1.0 \\
Debounced AND & 0.035 [0.027,\,0.044] & 0.56 [0.44,\,0.71] & 2.00 & 100 & 2.0 \\
2-of-4 rule & 0.020 [0.014,\,0.026] & 0.32 [0.23,\,0.43] & 1.00 & 100 & 3.9 \\
\textbf{Extension~I} & 0 ($<$0.025) & 0 ($<$0.025) & 2.00 & 100 & 4.0 \\
\midrule
\multicolumn{6}{l}{$q=0.02$} \\
Binary AND~\cite{soula2024iodt} & 7.8 [7.7,\,7.9] & 73 [72,\,73] & 1.00 & 100 & 1.0 \\
Debounced AND & 0.60 [0.56,\,0.63] & 8.7 [8.2,\,9.2] & 2.00 & 100 & 2.0 \\
2-of-4 rule & 0.25 [0.22,\,0.27] & 3.8 [3.5,\,4.2] & 1.00 & 100 & 3.7 \\
\textbf{Extension~I} & 0.014 [0.0078,\,0.020] & 0.18 [0.11,\,0.27] & 1.98 & 100 & 4.0 \\
\midrule
\multicolumn{6}{l}{$q=0.04$ (stress level)} \\
Binary AND~\cite{soula2024iodt} & 15 [15,\,15] & 93 [92,\,93] & 1.00 & 100 & 1.0 \\
Debounced AND & 2.1 [2.0,\,2.2] & 26 [25,\,27] & 2.00 & 100 & 2.0 \\
2-of-4 rule & 0.96 [0.92,\,1.0] & 14 [13,\,15] & 1.00 & 100 & 3.4 \\
\textbf{Extension~I} & 0.097 [0.079,\,0.12] & 1.1 [0.92,\,1.3] & 1.96 & 100 & 3.9 \\
\bottomrule
\end{tabular}
\end{table*}

Three regimes appear.  At reading~B the AND falsely isolates $0.39\%$
of honest nodes within 16 rounds, close to the $0.33\%$ that the
reading targets, and every rule that needs more than one failure
isolates none; there is little to improve.  At reading~A the AND
isolates $6.2\%$, while the debounced AND and the 2-of-4 rule isolate
$0.017\%$ and $0.0083\%$ (two and one of $12{,}000$ honest node-runs)
and Extension~I none.  At this level any rule that needs two failures
removes almost all false isolations, and the 2-of-4 rule does so
without delaying detection of a naive liar (a minimal-footprint liar
still takes it about four rounds).  Only as congestion grows does smoothing
make the difference: at $q=0.005$ Extension~I still records no false
flag where the debounced AND isolates $0.56\%$ and the 2-of-4 rule
$0.32\%$; at $q=0.02$ the figures are $0.18\%$, $8.7\%$ and $3.8\%$;
and at $q=0.04$ they are $1.1\%$, $26\%$ and $14\%$, against $93\%$
for the AND.

Fig.~\ref{fig:qsweep} shows why.  The fitted log--log slopes of the
per-round FPR are $0.99$ for the AND ($q$ from $0.0005$ to $0.02$),
$2.0$ for both two-failure rules ($q$ from $0.002$ to $0.04$), and
$3.1$ for Extension~I ($q$ from $0.02$ to $0.08$; below $0.02$ it
records too few flags to fit), in line with the $\Theta(q)$,
$\Theta(q^{2})$ and $\Theta(q^{3})$ orders of
Section~\ref{subsec:ext1-math}.  Equation~(\ref{eq:fpr_and}) predicts
$15.07\%$ for the AND at $q=0.04$, against $15.04\%$ measured.  The
price of smoothing is one round: Extension~I isolates a naive liar after
$1.96$--$2.00$ rounds, like the debounced AND, where the AND and the
2-of-4 rule need one.  Extension~I detects the minimal-footprint liar
in every run at every level, after $3.9$--$4.0$ rounds, as
Proposition~\ref{prop:persist} predicts.  Above $q\approx0.08$ every
aggregator degrades: at $q=0.12$ Extension~I falsely isolates $30\%$
of honest nodes within 16 rounds.

\Figure[!t](topskip=0pt,botskip=0pt,midskip=0pt)[width=\columnwidth]{figs/fig_qsweep.png}
{Per-round FPR (top, logarithmic axes) and false isolation within 16 rounds (bottom) against the per-factor congestion rate, with dotted lines at reading~A and at the stress level.\label{fig:qsweep}}

% ====================================================================
\subsection{Long Operation and Re-admission}
\label{sec:res-long}
A deployed trust engine runs for far longer than 16 rounds.
Table~\ref{tab:long} extends the horizon to
$4{,}096$ rounds ($600$ runs per cell, attackers honest).  False
isolation grows with the horizon for every rule: at $q=0.04$
Extension~I falsely isolates $1.2\%$ of honest nodes within 16 rounds
in these runs ($1.1\%$ in Table~\ref{tab:headline}), $20\%$ within
$256$ and $60\%$ within $1{,}024$, while the debounced AND and the
2-of-4 rule reach $99\%$ and $91\%$ within $256$ rounds.

The mean time to the first false isolation (MTTFI) does not depend on
the horizon.  At reading~A it is $250$ rounds for the AND (about
$29$~s at the assumed period), $6.2\times10^{4}$ rounds for the
debounced AND ($2.0$~h) and $1.5\times10^{5}$ for the 2-of-4 rule
($5.0$~h).  Extension~I recorded no flag in $1.5\times10^{7}$ honest
node-rounds, so its MTTFI is at least $4.9\times10^{6}$ rounds ($6.7$
days) at $95\%$ confidence.  At $q=0.005$ its MTTFI is
$7.4\times10^{5}$ rounds ($24$~h), against $2.6\times10^{3}$ rounds
($5$~min) and $6.6\times10^{3}$ rounds ($13$~min) for the two simple
rules; at $q=0.04$ it is $1.1\times10^{3}$ rounds ($2.2$~min), against
$52$ and $110$ rounds.

\begin{table}[!t]
\caption{Long operation (600 runs of 4{,}096 rounds; no attackers): false isolation within 1{,}024 rounds, MTTFI, and quarantine with re-admission.}
\label{tab:long}
\centering
\renewcommand{\arraystretch}{1.08}
\footnotesize
\setlength{\tabcolsep}{2.4pt}
\begin{tabular}{lrrrr}
\toprule
 & \textbf{False isol.} & \textbf{MTTFI} & \textbf{Quarant.} & \textbf{False quar.} \\
\textbf{Aggregator} & \textbf{1{,}024 rds \%} & \textbf{(rounds)} & \textbf{time \%} & \textbf{/node/h} \\
\midrule
\multicolumn{5}{l}{$q=0.001$} \\
AND~\cite{soula2024iodt} & 98 & 250 & 3.2 & 119 \\
Debounced AND & 1.7 & $6.2\times10^{4}$ & 0.013 & 0.49 \\
2-of-4 rule & 0.58 & $1.5\times10^{5}$ & 0.0052 & 0.20 \\
\textbf{Extension~I} & 0 & $\geq$$4.9\times10^{6}$ & 0 & 0 \\
\midrule
\multicolumn{5}{l}{$q=0.005$} \\
AND~\cite{soula2024iodt} & 100 & 51 & 15 & 520 \\
Debounced AND & 33 & 2{,}600 & 0.31 & 12 \\
2-of-4 rule & 15 & 6{,}600 & 0.12 & 4.7 \\
\textbf{Extension~I} & 0.14 & $7.4\times10^{5}$ & 0.0013 & 0.042 \\
\midrule
\multicolumn{5}{l}{$q=0.04$} \\
AND~\cite{soula2024iodt} & 100 & 7.8 & 73 & 1260 \\
Debounced AND & 100 & 52 & 15 & 524 \\
2-of-4 rule & 100 & 110 & 7.2 & 265 \\
\textbf{Extension~I} & 60 & 1{,}100 & 1.4 & 27 \\
\bottomrule
\end{tabular}
\end{table}

With quarantine and re-admission (Section~\ref{subsec:metrics}), a
falsely quarantined honest node returns after $8$ to $18$ rounds on
average ($9$ to $16$ under Extension~I).  At $q=0.04$ an honest node
then spends $1.4\%$ of its time in quarantine under Extension~I ($27$
false quarantines per node and hour), against $7.2\%$, $15\%$ and $73\%$
under the 2-of-4 rule, the debounced AND and the AND.  At $q=0.005$ the
share is $0.0013\%$ under Extension~I and $0.12$--$15\%$ under the
others.  Under sustained congestion as high as $q=0.04$, false
quarantines are an operating cost of every rule, which Extension~I cuts
by a factor of $5$ to $50$; at lower congestion it removes them.

% ====================================================================
\subsection{Stealthy, Minimal and Rule-Aware Attackers}
\label{sec:res-stealth}
Table~\ref{tab:stealth} varies the lying rate $p$ of attackers that lie
at random.  Stealth has a price for every rule except the AND, which
pays for its detection with false isolation.  At reading~A,
Extension~I detects a liar that fails $T$ and $W$ in $97\%$ of runs at
$p=0.5$ but in $57\%$ at $p=0.25$ and $13\%$ at $p=0.10$ within
16 rounds, much like the debounced AND; over 64 rounds the last two
figures rise to $97\%$ and $46\%$.  The 2-of-4 rule flags every lie
that fails two factors and detects such a liar almost as well as the
AND ($82\%$ of runs at $p=0.10$), at the false-isolation cost of
Table~\ref{tab:headline}.  A liar that fails only $W$ is
detected by Extension~I in $92\%$ of runs at $p=0.5$ and $32\%$ at
$p=0.25$ ($86\%$ over 64 rounds), about as often as by the 2-of-4
rule, whose persistence check is the same.  At $q=0.04$ every rule detects more, partly because
congestion adds failures to the attacker's own record.

\begin{table}[!t]
\caption{Detection within 16 rounds (\%) of attackers that lie at random with probability $p$ (2{,}000 runs; the false isolation of each rule is that of Table~\ref{tab:headline}).}
\label{tab:stealth}
\centering
\renewcommand{\arraystretch}{1.08}
\footnotesize
\setlength{\tabcolsep}{1.6pt}
\begin{tabular}{lrrrrrrrr}
\toprule
 & \multicolumn{4}{c}{\textbf{Lie fails $T$ and $W$}} & \multicolumn{4}{c}{\textbf{Lie fails $W$ only}} \\
\cmidrule(lr){2-5}\cmidrule(lr){6-9}
$p$ & AND & Deb. & 2-of-4 & Ext.~I & AND & Deb. & 2-of-4 & Ext.~I \\
\midrule
\multicolumn{9}{l}{$q=0.001$} \\
0.5 & $>$99.9 & 96 & $>$99.9 & 97 & $>$99.9 & 96 & 92 & 92 \\
0.25 & 99 & 55 & 99 & 57 & 99 & 55 & 33 & 32 \\
0.1 & 83 & 14 & 82 & 13 & 83 & 14 & 2.8 & 2.2 \\
\midrule
\multicolumn{9}{l}{$q=0.04$} \\
0.5 & 100 & 99 & $>$99.9 & 99 & 100 & 99 & 97 & 94 \\
0.25 & 99.9 & 82 & 99 & 70 & 99.9 & 82 & 64 & 44 \\
0.1 & 98 & 50 & 84 & 24 & 98 & 50 & 31 & 8.4 \\
\bottomrule
\end{tabular}
\end{table}

A rule-aware attacker does better (Table~\ref{tab:aware}).  The
simulated attacker predicts its own EMA from its past evidence
(assuming a $96\%$ PDR), keeps a margin of $0.02$ above $\theta$, and
respects every persistence window.  Without congestion it lied in $3$ of
$8$ rounds ($37.5\%$) against Extension~I and in $7$ of $32$ ($21.9\%$)
with a second window of $8$ failures in $32$ rounds, without being
flagged --- the budgets of Proposition~\ref{prop:persist}.  Every lie
still fails $W$, so the AND flags the first one, at a cost of $6.2\%$
false isolation at reading~A.  At reading~A congestion rarely adds
failures: over 256 rounds Extension~I detects the attacker in $37\%$ of
runs ($T$ and $W$) and $20\%$ ($W$ only), and with re-admission the
attacker delivers a lie in $36\%$ of rounds ($22\%$ with the second
window).  At $q=0.04$ congestion pushes the attacker over a trigger more
often, and the delivered share falls to $19\%$ ($14\%$); a rule-aware
liar that fails only $W$ delivers $25\%$ ($16\%$).  The simulated
attacker is tuned to Extension~I and faces every defence with the
same strategy, so the AND and debounced-AND figures show what that
attacker suffers, not what an attacker tuned to them would achieve;
against the debounced AND it delivered $16\%$ of rounds ($T$ and $W$)
and $4.5\%$ ($W$ only) at $q=0.04$.  Analytically, the simple rules
leave budgets of their own: the debounced AND never sees lies that
are one round apart (budget $50\%$), and the 2-of-4 rule flags every
lie that fails $T$ and $W$ but leaves a $W$-only liar the same
$37.5\%$ budget.  Threshold rules with published parameters thus leave
a lie budget that a second, longer window shrinks but does not remove.

\begin{table}[!t]
\caption{Rule-aware on-off attacker tuned to Extension~I (all defences face the same strategy): lie budget without congestion, detection within 256 rounds at reading~A (500 runs), lies delivered per round with re-admission (400 runs of 2{,}048 rounds, $T$ and $W$ fail), and the honest nodes' quarantine time at $q=0.04$.}
\label{tab:aware}
\centering
\renewcommand{\arraystretch}{1.08}
\footnotesize
\setlength{\tabcolsep}{2.0pt}
\begin{tabular}{lrrrrrrr}
\toprule
 & \multicolumn{2}{c}{\textbf{Budget \%}} & \multicolumn{2}{c}{\textbf{Det.\ \%}} & \multicolumn{2}{c}{\textbf{Delivered \%}} & \textbf{Hon.} \\
\cmidrule(lr){2-3}\cmidrule(lr){4-5}\cmidrule(lr){6-7}
\textbf{Defence} & $TW$ & $W$ & $TW$ & $W$ & $q{=}.001$ & $q{=}.04$ & \textbf{quar.\,\%} \\
\midrule
\textbf{Extension~I} & 37.5 & 37.5 & 37 & 20 & 36 & 19 & 1.4 \\
Ext.~I + 8 of 32 & 21.9 & 21.9 & 36 & 20 & 22 & 14 & 1.4 \\
AND~\cite{soula2024iodt} & 0 & 0 & 100 & 100 & 0.21 & 2.2 & 73 \\
\bottomrule
\end{tabular}
\end{table}

% ====================================================================
\subsection{Corrupted, Withholding and Slow Twins}
\label{sec:res-ext2}
Table~\ref{tab:collusion} turns to Extension~II at $q=0.04$, with
identical evidence for all replicas (B4).  A
corrupted single twin hides its node in every run.  With $3$-of-$5$
voting, up to $k-1=2$ corrupted replicas that shield the attackers and
frame honest nodes change nothing: detection stays at $100\%$ of runs
and of steady-state rounds, and the FPR and false isolation equal those
of an honest single twin, as Section~\ref{sec:ext2} argues.  A third
corrupted replica ($c=k$) defeats the vote: the attackers are never
flagged and every honest node is isolated.  Deployments that expect
three replicas at risk need $4$-of-$7$, which tolerates them.

Missing votes are handled by the quorum rule.  With two withholding
replicas plus $5\%$ vote timeouts, or with $20\%$ timeouts, it keeps
false isolation at the level of an honest single twin and detection at
$100\%$ of runs; the cost is rounds with no quorum ($14\%$ and
$5.8\%$), in which verdicts are held and an alarm is raised.  With
$20\%$ timeouts and two withholding replicas, half of all rounds have
no quorum and per-round detection falls to $97\%$, but no missing
vote caused a new isolation.  These timeout rows relax (B3); a node
that is already isolated stays isolated, and cannot be re-admitted,
while no quorum forms.

(B4) matters.  When each replica instead misreads each factor
independently with probability $\eta$ on top of the shared congestion
($q=0.04$), five honest replicas lower false isolation from $12\%$
(one twin) to $2.4\%$ at $\eta=0.05$, but replicas that frame honest
nodes now tip the nodes whose honest replicas disagree: false
isolation rises to $6.0\%$ with one framing replica and $29\%$ with
two, and to $72\%$ with two at $\eta=0.10$ (against $9.3\%$ with
five honest replicas).  Shielding costs little: with two shielding
replicas, steady-state detection falls from $99.9\%$ to $98\%$ at
$\eta=0.05$.  The ``honest twin'' guarantee of Extension~II therefore
holds only when the replicas see the same evidence.

\begin{table}[!t]
\caption{Corrupted, withholding and slow replicas at $q=0.04$, identical evidence (2{,}000 runs; per-round detection in steady state).}
\label{tab:collusion}
\centering
\renewcommand{\arraystretch}{1.08}
\footnotesize
\setlength{\tabcolsep}{2.0pt}
\begin{tabular}{p{0.40\columnwidth}rrrrr}
\toprule
 & \textbf{Det.} & \textbf{Per-rd} & \textbf{FPR} & \textbf{False} & \textbf{No} \\
\textbf{Setting} & \textbf{\%} & \textbf{det.\,\%} & \textbf{\%} & \textbf{isol.\,\%} & \textbf{quor.\,\%} \\
\midrule
Single twin, honest & 100 & 100 & 0.097 & 1.1 & -- \\
Single twin, corrupted & 0 & 0 & 0.097 & 1.1 & -- \\
3-of-5, $c=2$ shield + frame & 100 & 100 & 0.097 & 1.1 & 0 \\
3-of-5, $c=3$ shield + frame & 0 & 0 & 100 & 100 & 0 \\
4-of-7, $c=3$ shield + frame & 100 & 100 & 0.097 & 1.1 & 0 \\
\midrule
$c=2$ withhold, 5\% timeouts & 100 & 99.8 & 0.096 & 0.98 & 14 \\
20\% timeouts & 100 & $>$99.9 & 0.096 & 1.0 & 5.8 \\
$c=2$ withhold, 20\% timeouts & 100 & 97 & 0.091 & 0.67 & 49 \\
\bottomrule
\end{tabular}
\end{table}

% ====================================================================
\subsection{Audit Log: Cost and Tamper Evidence}
\label{sec:res-ext3}
Table~\ref{tab:chain} reports the audit log for $100$ sensors over
$30$ rounds of simulator output, with one corrupted replica and $5\%$
vote timeouts.  A header with three endorsements occupies $367$ to
$371$~bytes (DER signatures vary slightly in length) whatever the
number of nodes, so header storage falls from $11.0$~kB at $K=1$ to
$1.1$~kB at $K=10$; at the assumed period and $K=10$ this is about
$27$~MB a day for any fleet size.  Off-chain, the events, signed vote
vectors and state records of the 30 rounds occupy $115$--$152$~kB, of
which $87$~kB are events.

\begin{table}[!t]
\caption{Audit-log storage for 100 sensors over 30 rounds.}
\label{tab:chain}
\centering
\renewcommand{\arraystretch}{1.08}
\footnotesize
\setlength{\tabcolsep}{4pt}
\begin{tabular}{rrrrr}
\toprule
$K$ & \textbf{Headers} & \textbf{Header bytes} & \textbf{Off-chain (B)} & \textbf{Max header (B)} \\
\midrule
1 & 30 & 11{,}045 & 152{,}028 & 371 \\
5 & 6 & 2{,}205 & 120{,}834 & 370 \\
10 & 3 & 1{,}107 & 116{,}932 & 369 \\
20 & 2 & 734 & 115{,}620 & 367 \\
\bottomrule
\end{tabular}
\end{table}

Table~\ref{tab:tamper} reports $30$ random edits per cell at $K=10$.
The verifier rejected every edit when the adversary held fewer than
$k$ keys, whether it left the hashes untouched, recomputed them, or
re-signed with $k-1$ stolen keys.  Substituting interior-node values
for events and replaying a vote were rejected even with $k$ keys,
because record lengths and vote bindings are checked.  The other
$k$-key rows measure our implemented forger, not a limit of the
scheme: by Theorem~\ref{thm:tamper} a forger that re-signs every
affected record passes whenever no newer checkpoint exists.  Our
forger duplicated events and truncated windows without re-signing the
affected counts and votes, so it was rejected at $K=10$, whereas at
$K=1$ a truncation passed without a checkpoint.  It rewrote a trust
value, a state record or a vote and passed without a checkpoint or
with an older one, and was always rejected against a newer one.  Its
verdict edits were rejected without a checkpoint in $10\%$ of trials,
all of them rounds in which a stolen replica's vote had timed out and
the forger did not sign a replacement.  Finally, in $30$ trials a trust engine
that wrote a verdict its votes did not support never obtained the $k$
endorsements needed to commit the window.  These trials check the
implementation; the guarantee itself rests on the cryptographic
assumptions.

\begin{table}[!t]
\caption{Edits rejected by the verifier (\%, 30 per cell, $K=10$; the ``older'' column uses the 20--21 trials whose edit follows an earlier header) by edit and adversary; ``--'': not run.}
\label{tab:tamper}
\centering
\renewcommand{\arraystretch}{1.08}
\footnotesize
\setlength{\tabcolsep}{2.2pt}
\begin{tabular}{lrrrrrr}
\toprule
 & \multicolumn{2}{c}{\textbf{No keys}} & \textbf{$k{-}1$} & \multicolumn{3}{c}{\textbf{$k$ keys, checkpoint}} \\
\cmidrule(lr){2-3}\cmidrule(lr){5-7}
\textbf{Edit} & \textbf{edit} & \textbf{rehash} & \textbf{keys} & \textbf{none} & \textbf{older} & \textbf{newer} \\
\midrule
Verdict (with its votes) & 100 & 100 & 100 & 10 & 10 & 100 \\
Trust value $\tau_t$ & 100 & 100 & 100 & 0 & 0 & 100 \\
State record & 100 & 100 & 100 & 0 & 0 & 100 \\
Signed vote & 100 & 100 & 100 & 0 & 0 & 100 \\
Interior-node substitution & 100 & -- & -- & 100 & 100 & 100 \\
Duplicated event & 100 & -- & -- & 100 & 100 & 100 \\
Truncated window & 100 & -- & -- & 100 & 100 & 100 \\
Replayed vote & 100 & -- & -- & 100 & 100 & 100 \\
\bottomrule
\end{tabular}
\end{table}

% ====================================================================
\subsection{Robustness}
\label{sec:res-robust}
Fig.~\ref{fig:robust} relaxes assumptions (A1) and (A2) in four ways.

\emph{Bursts.}  Holding the mean congestion at $q=0.04$, a two-state
Markov chain alternates between $q=0.02$ and $q=0.20$ (one ninth of
the time in the high state) with mean bursts of $1$ to $32$ rounds.
Longer bursts raise Extension~I's false isolation from $3.7\%$ to
about $9\%$ when each node has its own chain (panel a), and from
$4.2\%$ to about $9.5\%$ when one chain governs the whole network.  The AND
($77$--$91\%$), the debounced AND ($18$--$29\%$) and the 2-of-4 rule
($16$--$32\%$) stay far higher in either case.

\emph{Common-mode events.}  An event that knocks out all factors of
every node behind a fog node, with probability $q_{\mathrm{corr}}$ per
round, cannot be told from an attack on all those nodes at once, and
Extension~I's false isolation rises to $8.4\%$ at
$q_{\mathrm{corr}}=0.005$ and $28\%$ at $0.02$ (panel b).  Because such
events hit many nodes together, we evaluated \emph{domain
suppression}: when more than half of a fog domain is flagged in the
same round, those flags are discarded for that round (an alarm would
be raised; it is not modelled).  It cuts
false isolation to $1.6\%$ and $3.5\%$ at the same rates and left the
detection of both attackers at $100\%$ of runs.  An adversary able to provoke
common-mode events could use suppression to mask its own lies during
them; that case is not evaluated.

\emph{Cross-factor events.}  When one node-level event fails $P$, $T$
and $W$ together ($4{,}000$ runs per point, no other noise), (A1) fails
across factors and Extension~I errs at $\Theta(q^{2})$ (panel c): its
FPR grows from $0.011\%$ to $0.055\%$ to $0.22\%$ as the event rate
doubles from $0.01$ to $0.04$.  The debounced AND is then better
at every rate ($0.15\%$ against $0.22\%$ at $0.04$), and the 2-of-4 rule (not plotted) is as bad
as the AND, since every event fails three factors.  Such an event also
looks exactly like a lie that fails $T$ and $W$ with a degraded PDR.

\emph{Uneven congestion.}  Concentrating the same mean congestion on
one factor ($0.13$ on it, $0.01$ on the others) leaves the AND
almost unchanged ($15.0\%$ to $15.6\%$ per round) but raises Extension~I's per-round FPR eight- to ninefold, to
$0.81$--$0.88\%$, and its false isolation to $4.8$--$5.7\%$ (panel d),
because one factor now fails repeatedly and the persistence check
fires.  Extension~I stays below the 2-of-4 rule ($11\%$), but a
deployment that knows its noisiest factor should set $m$ per factor.

\Figure[!t](topskip=0pt,botskip=0pt,midskip=0pt)[width=\textwidth]{figs/fig_robust.png}
{Robustness to bursty (a), common-mode (b), cross-factor (c) and uneven (d) congestion.\label{fig:robust}}

\emph{Scale.}  With $100$ sensors behind five fog nodes ($500$ runs, $q=0.04$),
Extension~I isolated $1.1\%$ of honest nodes, as in the eight-node
case, and with common-mode events at $q_{\mathrm{corr}}=0.005$ domain
suppression brought false isolation from $9.0\%$ back to $2.0\%$.

% ====================================================================
\subsection{Trace-Driven Evaluation on Public Datasets}
\label{sec:res-trace}
The congestion of the previous sections is independent across factors,
rounds and nodes.  To test the extensions on real loss patterns, we
replaced it, for every node, with the packet-delivery record of a real
deployment.  The Intel Lab trace~\cite{bodik2004intel} records which
readings of $54$ Mica2Dot motes reached a base station, one per mote
every $31$~s, over $36$ days of a lossy multi-hop network ($51\%$ of
readings arrived; timestamp jitter is not counted as loss, and periods
in which almost no mote was heard are logging gaps and are skipped).
The FlockLab dataset~\cite{jacob2019flocklab} records, about every two
hours in August 2019, which of $100$ back-to-back packets each of the
$27$ TelosB nodes of the FlockLab testbed~\cite{lim2013flocklab}
received from every other.  We took as fog gateways the two nodes that
put the most other nodes in range, and let each of the $24$ remaining
nodes in range report to the gateway it hears best (single hop, mean
packet reception $95\%$).  Two of these nodes ran only during the first
day and a half, so FlockLab windows use $22$ nodes and span the whole
month.

Each factor's evidence of a node in a round is one message, sent in up
to $r$ attempts that use consecutive transmission opportunities of
that node in the trace ($31$-s slots for Intel Lab, packets of one
test for FlockLab); the factor fails when all $r$ attempts are lost.
The retry count $r$ thus sets the loss level, which is then shaped
entirely by the trace.  The two traces differ in time scale: an Intel
Lab round lasts $4r$ slots ($8$ to $33$ minutes), whereas a FlockLab
round spans a few milliseconds of back-to-back packets and windows that
cross a test boundary join tests about two hours apart.  A round in
which no evidence of a node arrives at all is treated as
unavailability, and every rule keeps its previous verdict.  Each run
draws a window of $16$ rounds and $32$ (Intel Lab) or $22$ (FlockLab)
nodes alive throughout, two of which are attackers whose lies are added
to their own real losses ($1{,}000$ runs per cell, drawn from $73$ to
$309$ distinct windows).  No parameter was retuned.  Windows start on a
grid of $16$ rounds, and the intervals below resample whole windows.

\begin{table}[!t]
\caption{False isolation within 16 rounds (\%) on two public packet-delivery traces by retry count $r$, with $q$ the share of lost factor messages in the evaluated windows (1{,}000 trace runs; 2{,}000 runs for the i.i.d.\ column).}
\label{tab:trace}
\centering
\renewcommand{\arraystretch}{1.08}
\footnotesize
\setlength{\tabcolsep}{2.4pt}
\begin{tabular}{lrrrrrrrr}
\toprule
 & & & & & & & \multicolumn{2}{c}{\textbf{Ext.~I}} \\
\cmidrule(lr){8-9}
\textbf{Trace} & $r$ & $q$\,\% & AND & Deb. & 2-of-4 & Ext.~I & typ. & i.i.d. \\
\midrule
Intel Lab & 4 & 16 & 69 & 46 & 42 & \textbf{31} & 22 & 52 \\
 & 8 & 7.9 & 42 & 24 & 22 & \textbf{15} & 6.6 & 9.6 \\
 & 16 & 3.8 & 21 & 11 & 10 & \textbf{6.5} & 0.91 & 1.0 \\
\midrule
FlockLab & 1 & 5.5 & 48 & 22 & 30 & \textbf{8.3} & 0.67 & 3.2 \\
 & 2 & 3.2 & 35 & 8.4 & 6.8 & \textbf{2.7} & 0.074 & 0.60 \\
 & 4 & 1.9 & 8.0 & 2.6 & 3.1 & \textbf{2.0} & 0.15 & 0.16 \\
\bottomrule
\end{tabular}
\end{table}

Table~\ref{tab:trace} gives the result.  Extension~I has the lowest
false isolation in every cell, and the window-level interval of its
difference to each other rule excludes zero in every cell (the
smallest is $0.16$ to $1.1$ percentage points, against the debounced
AND at FlockLab with $r=4$).  The traces narrow its lead, however.  At
Intel Lab with $r=8$ ($7.9\%$ of messages lost) it falsely isolates
$15\%$ [$12$, $19$] of honest nodes within 16 rounds, against
$22$--$24\%$ for the two-failure rules and $42\%$ for the AND; with
$r=16$ the figures are $6.5\%$ [$4.2$, $9.1$], $10$--$11\%$ and $21\%$.
At FlockLab without retransmission it isolates $8.3\%$ [$7.5$, $9.2$]
against $22$--$30\%$ and $48\%$, and with $r=4$ $2.0\%$ [$1.4$, $2.7$]
against $2.6$--$3.1\%$ and $8.0\%$.  Independent congestion at the same
mean rate (column ``i.i.d.'') would make the AND isolate far more nodes
everywhere and Extension~I fewer everywhere except at the highest loss
level (Intel Lab, $r=4$).  Two properties of real loss explain this.
It is concentrated: at Intel Lab with $r=8$ the median node loses
$5.6\%$ of its messages over the whole trace and the lossiest tenth
$19\%$, so the AND's errors fall on a minority of nodes.  And it is
bursty across factors and rounds: in the evaluated windows two or more
factors of a node fail in the same round $1.4$ to $9.2$ times as often
as independent losses would make them, and a failed factor fails again
in the next round in $42$--$76\%$ of cases, which is the cross-factor
case of Section~\ref{sec:res-robust}, in which smoothing errs at
$\Theta(q^{2})$.  For the nodes whose long-run loss is at most the
median (column ``typ.'', about $10{,}000$ to $15{,}500$ node-runs per
cell), Extension~I's false isolation is $6.6\%$ and $0.91\%$ at Intel
Lab ($r=8$, $16$) and below $0.7\%$ at FlockLab; $65$--$99\%$ of its
false isolations fall on nodes with above-median loss.

Detection is slightly lower than in the synthetic study.  Extension~I
flagged a naive liar in $98.6$--$99.9\%$ of runs after about two rounds,
and a minimal-footprint liar in $98.8$--$99.9\%$ after $3.7$--$4.1$
rounds.  The misses are attacker nodes whose evidence was lost in
almost every round of the window, so that their verdict stayed at
\textsc{trust} under the silence rule; counting silent rounds as
failures instead detects every attacker and raises false isolation by
at most $2.4$ percentage points, with Extension~I still the lowest in
every cell.

Over longer windows false isolation grows for every rule: within $128$
rounds (Intel Lab, $r=8$) and $256$ rounds (FlockLab, $r=1$) Extension~I
isolates $30\%$ and $40\%$ of honest nodes against $41$--$66\%$ and
$61$--$81\%$ for the other rules, and with re-admission honest nodes
spend $14\%$ and $11\%$ of their time in quarantine against
$15$--$37\%$.  The rule-aware attacker, still tuned to Extension~I,
delivered a lie in $28\%$ (Intel Lab) and $27\%$ (FlockLab) of rounds
against it, $17\%$ with the second window, and was detected in
$61$--$74\%$ of runs; the AND and the 2-of-4 rule held it to about
$1.5\%$ of rounds at the cost of the false isolation above.  The
traces thus confirm the ordering of the synthetic study but not its
magnitude: on real, lossy links every rule isolates honest nodes far
more often, and per-node baselines that separate a weak link from
misbehaviour are needed before isolation is automated.

% ====================================================================
\subsection{Latency}
\label{sec:res-latency}
Table~\ref{tab:latency} gives the measured cryptographic cost per round.
Because each replica signs one vote vector per round, the engine
verifies five signatures, $2.0$~ms for $8$ or $100$ nodes, and each
replica signs once ($0.5$~ms); verifying one vote per node would cost the
engine $203$~ms per round at $100$ nodes, more than the assumed
$117$-ms period.  One EMA update for one node costs
$0.1$~$\mu$s; the persistence check was not timed separately.  The remaining term is the time
$d_{\mathrm{rep}}$ for votes to reach the engine, which depends on
replica placement and was not measured: replicas in different
availability zones of one cloud region are a single-digit number of
milliseconds apart~\cite{aws2024az}, whereas replicas at different fog
sites add a fog-to-cloud link ($100$~ms in the configuration
of~\cite{soula2024iodt}) unless the engine runs at the fog.
Extension~III adds nothing to the decision path.  Relative to the AND,
isolating a naive liar therefore takes one more round of evidence plus
about $2.5$~ms plus $d_{\mathrm{rep}}$.  All the added work runs in the
fog and cloud, so the extensions add no computation or energy on the
sensors.

\begin{table}[!t]
\caption{Measured per-round cost on the decision path (ms, $n=5$ replicas, single thread).}
\label{tab:latency}
\centering
\renewcommand{\arraystretch}{1.08}
\footnotesize
\setlength{\tabcolsep}{3pt}
\begin{tabular}{lrr}
\toprule
\textbf{Operation} & \textbf{8 nodes} & \textbf{100 nodes} \\
\midrule
EMA update, one node & \multicolumn{2}{c}{0.0001} \\
Replica: sign one vote vector & 0.52 & 0.48 \\
Engine: verify $5$ vote vectors & 2.0 & 2.0 \\
\bottomrule
\end{tabular}
\end{table}

% ====================================================================
\section{Discussion}
\label{sec:discussion}
% ====================================================================

% ====================================================================
\subsection{Comparison With Related Frameworks}
\label{sec:res-comparison}
Table~\ref{tab:comparison} positions the proposed system against the
trust frameworks closest to it.  Each cell reflects what the cited
paper describes, and the latency row reports what each paper measures.
Temporal smoothing is marked $\checkmark$ for an explicit forgetting
factor and $\circ$ for trust that accumulates history without one.
The latency figures are heterogeneous (trust set-up, block operations,
detect-and-isolate time) and come from different platforms, so they
show where each system places its costs rather than ranking the
systems.  Within this set, the proposed design is the only one that
combines a continuous, smoothed trust score, replicated twins with
voting, and a tamper-evident record kept off the decision path.

\begin{table*}[!t]
\caption{Comparison with related trust frameworks ($\checkmark$ yes, $\circ$ partial or optional, -- no).}
\label{tab:comparison}
\centering
\renewcommand{\arraystretch}{1.12}
\setlength{\tabcolsep}{6pt}
\footnotesize
\begin{tabular}{lccccccc}
\toprule
\textbf{Capability} & \textbf{BPT}~\cite{dewanta2021bpt} & \textbf{Debe}~\cite{debe2019iot}
 & \textbf{Hussain}~\cite{hussain2020context} & \textbf{Onwubiko}~\cite{onwubiko2023enabling}
 & \textbf{TMIoDT}~\cite{akram2024tmiodt} & \textbf{IoDT}~\cite{soula2024iodt} & \textbf{Proposed} \\
\midrule
Digital twin              & -- & -- & -- & \checkmark & \checkmark & \checkmark & \checkmark \\
Continuous trust score    & \checkmark & \checkmark & \checkmark & -- & $\circ$ & -- & \checkmark \\
Temporal smoothing        & $\circ$ & $\circ$ & \checkmark & -- & $\circ$ & -- & \checkmark \\
Replicated twins, voting  & -- & -- & -- & -- & -- & -- & \checkmark \\
Blockchain                & \checkmark & \checkmark & -- & \checkmark & \checkmark & -- & $\circ$ \\
Chain off decision path   & $\circ$ & -- & n/a & -- & -- & n/a & \checkmark \\
Tamper-evident record     & \checkmark & \checkmark & -- & \checkmark & \checkmark & -- & \checkmark \\
\midrule
Reported latency          & 2.47\,s & n.r. & n.r. & n.r. & 15--16\,ms & 250--300\,ms & +1 round \\
(what it measures)        & trust set-up & & & (gas cost only) & block op. & detect + isolate & +2.5\,ms\,+\,$d_{\mathrm{rep}}$ \\
\bottomrule
\multicolumn{8}{l}{n.r.\ = not reported.} \\
\end{tabular}
\end{table*}

TMIoDT~\cite{akram2024tmiodt} is the closest precedent.  It assigns
each drone a blockchain-recorded credit score and validates
transactions by proof-of-authority consensus, so the chain is involved
in its trust decisions, and it detects intrusions with deep models trained
on labelled traffic (about $99.1\%$ accuracy on ToN-IoT); its twins are
not replicated.  The proposed score needs no training data and keeps
the log off the decision path, but sees only the four behavioural
factors, so the two approaches are complementary.

\subsection{Limitations}
Behaviour-based trust has hard limits, which the experiments make
explicit.  An attacker that knows the rule keeps a lie budget
($37.5\%$ of rounds against Extension~I without congestion,
$21.9\%$ with a second window), low-rate random liars are detected
slowly, and one that never deviates on any factor cannot be seen at
this layer at all.  Bad-mouthing through the notification factor is
not covered.  Common-mode events need domain suppression, which an
adversary able to provoke them could abuse, and events that fail
several factors at once lower the order of the error to
$\Theta(q^{2})$.  On the architecture side, the fog relay is trusted
(it enforces verdicts and logs the verdicts it enforced), although a
replica hosted at a fog site may be corrupted, and the trust engine is
trusted for timeliness; its integrity is checked after the fact,
within $K$ rounds.  Extension~II assumes identical evidence and
bounded vote delay; with independent replica errors, two framing
replicas raise false isolation several-fold.  An adversary that controls
$k$ replicas defeats any quorum by construction.  Finally, the
evaluation is a simulation of the trust logic; honest losses come from
synthetic congestion or from two public delivery traces mapped onto the
factors by our retransmission model, attackers are simulated, and on
those traces lossy nodes are falsely isolated often under every rule;
under the silence rule a node that goes completely silent keeps its
last verdict, so silence itself must be monitored separately.
The round length, the header append
latency and the replica delay $d_{\mathrm{rep}}$ are assumptions or
were not measured, and no ledger was deployed.
% ====================================================================
\section{Conclusion}
\label{sec:conclusion}
% ====================================================================
The IoDT model of~\cite{soula2024iodt} pairs every IoT node with a
cloud-resident twin and decides trust from four interpretable factors.
Its binary AND, single twin and mutable verdict store make it fragile
under congestion, blind to a corrupted twin, and not tamper-evident.
This paper kept the factor set and added three extensions: a weighted,
EMA-smoothed score with a per-factor persistence check; quorum voting
over signed vote vectors of replicated twins; and a signed Merkle
audit log of verdicts and votes, endorsed by the replicas off the
decision path.

The evaluation compared the score with delay-matched baselines rather
than with the AND alone, and the comparison changes the claim.  At the
congestion levels implied by the original evaluation, any rule that
needs two failures removes almost all false isolations, and a
2-of-4 rule does so without delaying the detection of a naive liar.  Smoothing pays off as
congestion grows and over long operation: at $4\%$ per-factor
congestion Extension~I falsely isolated $1.1\%$ of honest nodes within
16 rounds, against $14$--$26\%$ for the simple rules and $93\%$ for the
AND, its mean time to a first false isolation was $110$ times longer
than that of the better simple rule at $0.5\%$ congestion and $10$
times longer at $4\%$, and with re-admission honest nodes spent
$1.4\%$ of their time in quarantine at $4\%$,
at the cost of one extra round of evidence.  An operator who knows
that congestion is low can therefore use the simpler 2-of-4 rule; one
who does not, or who expects congestion to vary, gains from
smoothing.  On two public packet-delivery traces, Extension~I kept the
lowest false isolation in every setting, but real losses, bursty and
concentrated on a minority of nodes, raised it to $2.0$--$31\%$ within
16 rounds, mostly on nodes with above-median loss.  When the replicas of a twin see identical evidence, a quorum rule
kept detection at the level of an honest twin with up to two of five
replicas corrupted, and missing or withheld votes caused no new
isolation; with independent replica errors, framing replicas raise
false isolation markedly.  The audit log rejected every edit by an
adversary with fewer than $k$ replica keys, rejected interior-node
substitutions and replayed votes even with $k$ keys, and exposed
unsupported verdicts; its header size is independent of the number of
nodes.

The results are bounded by the limitations listed in
Section~\ref{sec:discussion}: a rule-aware attacker keeps a lie budget
of $37.5\%$ of rounds without congestion ($21.9\%$ with a second
window), bad-mouthing through the notification factor is not covered,
common-mode events need domain suppression, cross-factor events lower
the order of the error from $\Theta(q^{3})$ to $\Theta(q^{2})$,
voting relies on identical evidence, the fog relay is trusted, lossy
nodes are falsely isolated often on real traces, and the evaluation
simulates the trust logic with assumed timing and no deployed ledger.  Future work will
(i)~replace fixed thresholds with sequential tests such as
CUSUM~\cite{page1954continuous} whose parameters are randomised, to
shrink the lie budget; (ii)~corroborate notification failures across
several destinations against bad-mouthing; (iii)~calibrate thresholds to each node's own loss history so that
a weak link is not taken for misbehaviour; (iv)~stack the behavioural
score with a traffic-level detector such
as~\cite{meidan2018nbaiot}; (v)~deploy the replicas and the log on a
multi-domain fog testbed to measure $d_{\mathrm{rep}}$ and the append
latency; and (vi)~couple the audit log with a SOAR
platform~\cite{empl2022soar4iot} for automated response.

\section*{Acknowledgment}
The author thanks P.~Bodik, W.~Hong, C.~Guestrin, S.~Madden,
M.~Paskin and R.~Thibaux for the Intel Lab data, and R.~Jacob,
R.~Da~Forno, R.~Tr{\"u}b, A.~Biri and L.~Thiele for the FlockLab
link-quality dataset.

\section*{Data and Code Availability}
The simulator, the experiment, table and figure scripts, the
configuration of every run and all raw result files are provided as
supplementary material with this submission and will be archived in a
public repository with a persistent identifier.  The trace-driven
experiments use the Intel Lab data~\cite{bodik2004intel}, reused with
the acknowledgement its authors request (above), and the FlockLab link-quality
dataset~\cite{jacob2019flocklab} (CC~BY~4.0); the pre-processing scripts
and the exact files used are included.

\appendices
\section{Proofs}
\label{app:proofs}

\subsection{False-positive rate of the EMA score}
This subsection proves Proposition~\ref{prop:fpr_ema}.
Equation (\ref{eq:sstar}) follows by solving $\tau_t<\theta$ for $s_t$
in (\ref{eq:ema}).  Unrolling (\ref{eq:ema}) over the last $L$ rounds
gives
\[
\tau_t=\alpha\sum_{j=0}^{L-1}(1-\alpha)^{j}s_{t-j}+(1-\alpha)^{L}\tau_{t-L},
\qquad \tau_{t-L}\ge 0 .
\]
With $\underline P=92$ a clean round scores at least $0.976$, a round
with one failure at least $0.676$ (energy), and a round with two
failures at least $0.476$ ($E$ with $T$, or $E$ with $W$).  If at most
two failures occur in rounds $t-L+1,\dots,t$, the sum is smallest when
both fall in round $t$, and with $\alpha=0.5$ and $L=6$
\[
\tau_t\ \ge\ 0.5\cdot0.476+0.976\sum_{j=1}^{5}0.5^{\,j+1}
\ =\ 0.711\ >\ \theta=0.70,
\]
whatever happened before round $t-5$.  A flag therefore needs at
least three failures among the $4L=24$ factor-rounds of the last six
rounds, which under (A1)--(A2) has probability at most
$\binom{24}{3}q^{3}$.  Conversely, $E$, $T$ and $W$ failing together
in round $t$ give $s_t\le0.30$ and $\tau_t\le0.5\cdot0.30+0.5=0.65$,
which flags the node, so the rate is also $\Omega(q^{3})$.
For other parameters, $E$ and $T$ failing together in round $t$ give
$s_t\le0.5$ and $\tau_t\le0.5\,\alpha+(1-\alpha)\,\tau_{t-1}$, where
$\tau_{t-1}<1$ whenever some recent PDR was below $100\%$.  At
$\theta=0.70$ and $\alpha=0.6$ this gives $\tau_t\le0.3+0.4\,\tau_{t-1}<0.70$,
and at $\alpha=0.5$ and $\theta=0.75$ it gives
$\tau_t\le0.25+0.5\,\tau_{t-1}<0.75$ (larger $\alpha$ or $\theta$ only
lowers the margin); in both cases two failures in one round flag the
node and the order is $\Theta(q^{2})$.

\subsection{Persistence check}
This subsection proves Proposition~\ref{prop:persist}.
For one factor, $m$ failures among $w$ independent rounds have
probability $\Pr[\mathrm{Bin}(w,q)\ge m]\le\binom{w}{m}q^{m}$; the union
bound over four factors gives the first claim.  The second follows
because each lie is a failure of that factor.  The check fires only
when some window holds $m$ failures of one factor, so a liar that
keeps at most $m-1$ lies in every window, and meets no congestion
failure of the same factor, is never flagged by it; spacing lies
evenly attains the rate $(m-1)/w$, and with two windows the smaller
of the two budgets applies.

\subsection{Tamper evidence}
This subsection proves Theorem~\ref{thm:tamper}.
The header commits to the counts and to RFC~6962 Merkle roots, whose
leaf and interior hashes are domain-separated and whose shape is fixed
by the leaf count; two different sequences of the same length
therefore have different roots unless a hash collision occurs, and a
change of length changes the committed count.  Every record has a fixed
length that $V$ checks, so interior-node values cannot stand in for
records.  Any modification, insertion, deletion or reordering in $B_i$
thus changes a root or a count, and hence the header, unless a
collision occurs.  The original endorsements cover the old header, so
a header that is not re-signed fails verification; re-signing needs
$k$ valid signatures, which an adversary with fewer than $k$ keys can
produce only by forgery.  A modified vote likewise needs the signature
of its replica, and a modified verdict that its votes do not support is
rejected by the quorum check.  An adversary with $k$ keys can re-sign
the votes of its replicas and the headers from $B_i$ onward and recompute
all hashes, producing an internally consistent log; then $h_j$ changes
for every $j\geq i$ and differs from a previously published
$h^{\star}_j$ unless a collision occurs, while a checkpoint with
$j<i$ is unaffected.


% ====================================================================
\begin{thebibliography}{99}

\bibitem{neshenko2019demystifying}
N.~Neshenko, E.~Bou-Harb, J.~Crichigno, G.~Kaddoum, and N.~Ghani,
``Demystifying IoT security: An exhaustive survey on IoT vulnerabilities
and a first empirical look on Internet-scale IoT exploitations,''
\textit{IEEE Commun. Surveys Tuts.}, vol.~21, no.~3,
pp.~2702--2733, 3rd Quart., 2019.

\bibitem{tao2019digital}
F.~Tao, H.~Zhang, A.~Liu, and A.~Y.~C. Nee,
``Digital twin in industry: State-of-the-art,''
\textit{IEEE Trans. Ind. Informat.}, vol.~15, no.~4,
pp.~2405--2415, Apr.~2019.

\bibitem{minerva2020digital}
R.~Minerva, G.~M. Lee, and N.~Crespi,
``Digital twin in the IoT context: A survey on technical features,
scenarios, and architectural models,''
\textit{Proc. IEEE}, vol.~108, no.~10, pp.~1785--1824, Oct.~2020.

\bibitem{soula2024iodt}
M.~Soula, B.~Mbarek, and A.~Meddeb,
``A real-time trust management model using digital twin in IoT networks,''
\textit{IEEE Access}, vol.~12, pp.~183326--183343, 2024,
doi: 10.1109/ACCESS.2024.3506870.

\bibitem{bao2012dynamic}
F.~Bao and I.-R. Chen,
``Dynamic trust management for Internet of Things applications,''
in \textit{Proc. Int. Workshop Self-Aware Internet of Things (Self-IoT)},
ACM, 2012, pp.~1--6.

\bibitem{nitti2014trustworthiness}
M.~Nitti, R.~Girau, and L.~Atzori,
``Trustworthiness management in the Social Internet of Things,''
\textit{IEEE Trans. Knowl. Data Eng.}, vol.~26, no.~5,
pp.~1253--1266, May~2014.

\bibitem{sun2006trust} Y.~L.~Sun, Z.~Han, W.~Yu, and K.~J.~R.~Liu, ``A trust evaluation framework in distributed networks: Vulnerability analysis and defense against attacks,'' in \emph{Proc. IEEE INFOCOM 2006, 25th IEEE Int. Conf. Comput. Commun.}, Apr.~2006, pp.~1--13, doi: 10.1109/INFOCOM.2006.154.

\bibitem{alcaraz2022digital}
C.~Alcaraz and J.~Lopez,
``Digital twin: A comprehensive survey of security threats,''
\textit{IEEE Commun. Surveys Tuts.}, vol.~24, no.~3,
pp.~1475--1503, 3rd Quart., 2022.

\bibitem{dorri2017blockchain}
A.~Dorri, S.~S. Kanhere, R.~Jurdak, and P.~Gauravaram,
``Blockchain for IoT security and privacy: The case study of a smart home,''
in \textit{Proc. IEEE Int. Conf. Pervasive Comput. Commun. Workshops
(PerCom Workshops)}, 2017, pp.~618--623.

\bibitem{novo2018blockchain}
O.~Novo,
``Blockchain meets IoT: An architecture for scalable access management
in IoT,''
\textit{IEEE Internet Things J.}, vol.~5, no.~2,
pp.~1184--1195, Apr.~2018.

\bibitem{kumar2020blockedge}
T.~Kumar \textit{et al.},
``BlockEdge: Blockchain-edge framework for industrial IoT networks,''
\textit{IEEE Access}, vol.~8, pp.~154166--154185, 2020,
doi: 10.1109/ACCESS.2020.3017891.

\bibitem{dewanta2021bpt}
F.~Dewanta and M.~Mambo,
``BPT scheme: Establishing trusted vehicular fog computing service
for rural area based on blockchain approach,''
\textit{IEEE Trans. Veh. Technol.}, vol.~70, no.~2,
pp.~1752--1769, Feb.~2021,
doi: 10.1109/TVT.2021.3051258.

\bibitem{akram2024tmiodt}
J.~Akram, A.~Anaissi, R.~S. Rathore, R.~H. Jhaveri, and A.~Akram,
``Digital twin-driven trust management in Open RAN-based spatial
crowdsourcing drone services,''
\textit{IEEE Trans. Green Commun. Netw.}, vol.~8, no.~3,
pp.~1061--1075, Sep.~2024,
doi: 10.1109/TGCN.2024.3403901.

\bibitem{sleuters2019digital}
J.~Sleuters, Y.~Li, J.~Verriet, M.~Velikova, and R.~Doornbos,
``A digital twin method for automated behavior analysis of
large-scale distributed IoT systems,''
in \textit{Proc. 14th Annu. Conf. Syst. Syst. Eng. (SoSE)},
2019, pp.~7--12, doi: 10.1109/SYSOSE.2019.8753845.

\bibitem{gehrmann2020digital}
C.~Gehrmann and M.~Gunnarsson,
``A digital twin based industrial automation and control system
security architecture,''
\textit{IEEE Trans. Ind. Informat.}, vol.~16, no.~1,
pp.~669--680, Jan.~2020.

\bibitem{empl2022soar4iot}
P.~Empl, D.~Schlette, D.~Zupfer, and G.~Pernul,
``SOAR4IoT: Securing IoT assets with digital twins,''
in \textit{Proc. 17th Int. Conf. Availability, Rel. Secur. (ARES)},
ACM, 2022, pp.~1--10, doi: 10.1145/3538969.3538975.

\bibitem{onwubiko2023enabling} A.~Onwubiko, R.~Singh, S.~Awan, Z.~Pervez, and N.~Ramzan, ``Enabling trust and security in digital twin management: A blockchain-based approach with Ethereum and IPFS,'' \emph{Sensors}, vol.~23, no.~14, Art.~no.~6641, Jul.~2023, doi: 10.3390/s23146641.

\bibitem{sasikumar2023blockchain} A.~Sasikumar \emph{et al.}, ``Blockchain-based trust mechanism for digital twin empowered Industrial Internet of Things,'' \emph{Future Gener. Comput. Syst.}, vol.~141, pp.~16--27, Apr.~2023, doi: 10.1016/j.future.2022.11.002.

\bibitem{kamvar2003eigentrust}
S.~D. Kamvar, M.~T. Schlosser, and H.~Garcia-Molina,
``The EigenTrust algorithm for reputation management in P2P networks,''
in \textit{Proc. 12th Int. Conf. World Wide Web (WWW)},
ACM, 2003, pp.~640--651.

\bibitem{josang2002beta} A.~J{\o}sang and R.~Ismail, ``The beta reputation system,'' in \emph{Proc. 15th Bled Electron. Commerce Conf.}, Bled, Slovenia, 2002, Art.~no.~41. [Online]. Available: \url{https://aisel.aisnet.org/bled2002/41}

\bibitem{chen2016trust} I.-R.~Chen, F.~Bao, and J.~Guo, ``Trust-based service management for social Internet of Things systems,'' \emph{IEEE Trans. Dependable Secure Comput.}, vol.~13, no.~6, pp.~684--696, 2016, doi: 10.1109/TDSC.2015.2420552.

\bibitem{yan2014survey}
Z.~Yan, P.~Zhang, and A.~V. Vasilakos,
``A survey on trust management for Internet of Things,''
\textit{J. Netw. Comput. Appl.}, vol.~42, pp.~120--134, 2014,
doi: 10.1016/j.jnca.2014.01.014.

\bibitem{guo2017survey} J.~Guo, I.-R.~Chen, and J.~J.~P.~Tsai, ``A survey of trust computation models for service management in Internet of Things systems,'' \emph{Comput. Commun.}, vol.~97, pp.~1--14, 2017, doi: 10.1016/j.comcom.2016.10.012.

\bibitem{konsta2023survey}
A.~M. Konsta, A.~Lluch Lafuente, and N.~Dragoni,
``A survey of trust management for Internet of Things,''
\textit{IEEE Access}, vol.~11, pp.~122175--122204, 2023,
doi: 10.1109/ACCESS.2023.3327335.

\bibitem{hussain2020context}
Y.~Hussain \textit{et al.},
``Context-aware trust and reputation model for fog-based IoT,''
\textit{IEEE Access}, vol.~8, pp.~31622--31632, 2020,
doi: 10.1109/ACCESS.2020.2972968.

\bibitem{bonomi2012fog}
F.~Bonomi, R.~Milito, J.~Zhu, and S.~Addepalli,
``Fog computing and its role in the Internet of Things,''
in \textit{Proc. 1st ACM Workshop Mobile Cloud Comput. (MCC)},
2012, pp.~13--16.

\bibitem{chiang2016fog}
M.~Chiang and T.~Zhang,
``Fog and IoT: An overview of research opportunities,''
\textit{IEEE Internet Things J.}, vol.~3, no.~6,
pp.~854--864, Dec.~2016.

\bibitem{shi2016edge}
W.~Shi, J.~Cao, Q.~Zhang, Y.~Li, and L.~Xu,
``Edge computing: Vision and challenges,''
\textit{IEEE Internet Things J.}, vol.~3, no.~5,
pp.~637--646, Oct.~2016.

\bibitem{debe2019iot}
M.~Debe, K.~Salah, M.~H.~U. Rehman, and D.~Svetinovic,
``IoT public fog nodes reputation system: A decentralized solution
using Ethereum blockchain,''
\textit{IEEE Access}, vol.~7, pp.~178082--178093, 2019,
doi: 10.1109/ACCESS.2019.2958355.

\bibitem{haber1991time} S.~Haber and W.~S.~Stornetta, ``How to time-stamp a digital document,'' \emph{J. Cryptol.}, vol.~3, no.~2, pp.~99--111, 1991, doi: 10.1007/BF00196791.

\bibitem{schneier1999secure} B.~Schneier and J.~Kelsey, ``Secure audit logs to support computer forensics,'' \emph{ACM Trans. Inf. Syst. Secur.}, vol.~2, no.~2, pp.~159--176, May~1999, doi: 10.1145/317087.317089.

\bibitem{crosby2009efficient} S.~A.~Crosby and D.~S.~Wallach, ``Efficient data structures for tamper-evident logging,'' in \emph{Proc. 18th USENIX Secur. Symp.}, Montreal, QC, Canada, Aug.~2009. [Online]. Available: \url{https://www.usenix.org/conference/usenixsecurity09/technical-sessions/presentation/efficient-data-structures-tamper-evident}

\bibitem{laurie2013rfc6962} B.~Laurie, A.~Langley, and E.~Kasper, ``Certificate Transparency,'' Internet Engineering Task Force, RFC~6962, Jun.~2013, doi: 10.17487/RFC6962. [Online]. Available: \url{https://www.rfc-editor.org/rfc/rfc6962}

\bibitem{buldas2013ksi}
A.~Buldas, A.~Kroonmaa, and R.~Laanoja,
``Keyless signatures' infrastructure: How to build global
distributed hash-trees,''
in \textit{Proc.\ Nordic Conf.\ Secure IT Systems (NordSec)},
ser. LNCS, vol.~8208, Springer, 2013, pp.~313--320,
doi: 10.1007/978-3-642-41488-6\_21.

\bibitem{kalodner2018arbitrum}
H.~Kalodner, S.~Goldfeder, X.~Chen, S.~M. Weinberg, and E.~W. Felten,
``Arbitrum: Scalable, private smart contracts,''
in \textit{Proc.\ 27th USENIX Security Symposium},
2018, pp.~1353--1370.

\bibitem{meidan2018nbaiot}
Y.~Meidan \textit{et al.},
``N-BaIoT: Network-based detection of IoT botnet attacks using deep
autoencoders,''
\textit{IEEE Pervasive Comput.}, vol.~17, no.~3, pp.~12--22, Jul.--Sep.~2018,
doi: 10.1109/MPRV.2018.03367731.

\bibitem{douceur2002sybil} J.~R.~Douceur, ``The Sybil attack,'' in \emph{Peer-to-Peer Systems (IPTPS 2002)}, ser. Lecture Notes in Computer Science, vol.~2429. Berlin, Germany: Springer, 2002, pp.~251--260, doi: 10.1007/3-540-45748-8\_24.

\bibitem{friedman2001social} E.~J.~Friedman and P.~Resnick, ``The social cost of cheap pseudonyms,'' \emph{J. Econ. Manage. Strategy}, vol.~10, no.~2, pp.~173--199, Jun.~2001, doi: 10.1111/j.1430-9134.2001.00173.x.

\bibitem{malkhi1998byzantine} D.~Malkhi and M.~K.~Reiter, ``Byzantine quorum systems,'' \emph{Distrib. Comput.}, vol.~11, no.~4, pp.~203--213, Oct.~1998, doi: 10.1007/s004460050050.

\bibitem{castro2002practical}
M.~Castro and B.~Liskov,
``Practical Byzantine fault tolerance and proactive recovery,''
\textit{ACM Trans. Comput. Syst.}, vol.~20, no.~4,
pp.~398--461, Nov.~2002.

\bibitem{gupta2017ifogsim}
H.~Gupta, A.~V. Dastjerdi, S.~K. Ghosh, and R.~Buyya,
``iFogSim: A toolkit for modeling and simulation of resource
management techniques in the Internet of Things, edge and fog
computing environments,''
\textit{Softw., Pract. Exper.}, vol.~47, no.~9,
pp.~1275--1296, 2017.

\bibitem{bodik2004intel} P.~Bodik, W.~Hong, C.~Guestrin, S.~Madden, M.~Paskin, and R.~Thibaux, ``Intel Lab Data,'' Intel Berkeley Research Lab, 2004. Accessed: Oct.~2026. [Online]. Available: \url{https://db.csail.mit.edu/labdata/labdata.html}

\bibitem{jacob2019flocklab} R.~Jacob, R.~Da~Forno, R.~Tr{\"u}b, A.~Biri, and L.~Thiele, ``Dataset: Wireless link quality estimation on FlockLab -- and beyond,'' in \emph{Proc. 2nd Workshop Data Acquisition To Analysis (DATA)}, Nov.~2019, pp.~57--60, doi: 10.1145/3359427.3361907. Data: doi: 10.5281/zenodo.3731498.

\bibitem{lim2013flocklab} R.~Lim, F.~Ferrari, M.~Zimmerling, C.~Walser, P.~Sommer, and J.~Beutel, ``FlockLab: A testbed for distributed, synchronized tracing and profiling of wireless embedded systems,'' in \emph{Proc. 12th Int. Conf. Inf. Process. Sensor Netw. (IPSN)}, Apr.~2013, pp.~153--166, doi: 10.1145/2461381.2461402.

\bibitem{aws2024az} Amazon Web Services, ``Availability Zones,'' in \emph{AWS Fault Isolation Boundaries}, AWS Whitepaper. Accessed: Oct.~2026. [Online]. Available: \url{https://docs.aws.amazon.com/whitepapers/latest/aws-fault-isolation-boundaries/availability-zones.html}

\bibitem{page1954continuous} E.~S.~Page, ``Continuous inspection schemes,'' \emph{Biometrika}, vol.~41, no.~1/2, pp.~100--115, 1954, doi: 10.1093/biomet/41.1-2.100.

\end{thebibliography}

\begin{IEEEbiography}[{\includegraphics[width=1in,height=1.25in,clip,keepaspectratio]{author_alwabel.jpg}}]{ABDULELAH ALWABEL}
(Member, IEEE) is an Associate Professor of Cloud Computing and
Digital Transformation with the Department of Computer Sciences,
College of Computer Engineering and Sciences, Prince Sattam bin
Abdulaziz University, Al-Kharj, Saudi Arabia.  His research interests
include cloud, fog and edge computing, the Internet of Things, digital
twins, trust management, and the application of artificial
intelligence to distributed systems.  He supervises M.Sc.\ and Ph.D.\
students in these areas, and his work also addresses enterprise IT
and digital transformation in the public sector.
\end{IEEEbiography}

\EOD

\end{document}
